
\section{\BookRecordIndex{T0324}图像伪循环}\label{sec:6-7}

上一节的结果不能用于\BookRecordIndex{T0058}半正流形中的所有类 $A$，主要原因是主\BookRecordIndex{T0246}模空间
$\mathcal M_{0,k}(A;J)$ 中可能含有非单纯元素。这种情形只会在 $A$ 本身是倍数类时发生。
处理这个问题有若干可能的方法。其思路是修改基本方程 $\bar\partial_J(u)=0$，使它不再有
多重覆盖解：或者加入一个非齐次扰动项，或者让 $J$ 随定义域中的点 $z$ 变化。本节采用后
一种方法。这是第~8 章所述更一般构造的一个特例；那个构造还涉及一些在几何上颇有意思的
扰动。

我们将尽可能对任意亏格的定义域 $(\Sigma,j_\Sigma)$ 展开论证。不过，本节的主要结果
（定理~\ref{thm:6-7-1}）只对亏格零成立。原因在于，高亏格常值曲线无法通过\BookRecordIndex{T0088}殆复结构的
通有（依赖于 $z$ 的）扰动达到横截；如第~8 章所述，这需要更一般的扰动。映入二维球面的
次数一亏格一曲线清楚地说明了这一现象：这样的曲线并不存在，但用 Hamilton 扰动定义的
\BookRecordIndex{T0143}格罗莫夫--威滕不变量却非零（例 8.6.12）。

\paragraph{\BookRecordIndex{T0324}图像伪循环。}
本节研究的基本对象，是方程
\begin{equation}\label{eq:6-7-1}
  du(z)+J_z(u(z))\circ du(z)\circ j_\Sigma(z)=0
\end{equation}
的解 $u:\Sigma\to M$ 所成的空间，其中 $z\in\Sigma$，$j=j_\Sigma$ 是 $\Sigma$ 上固定的
复结构，而
\[
  \mathbf J=\{J_z\}_{z\in\Sigma}
\]
是由 $\Sigma$ 中的点参数化的、$M$ 上被 $\omega$ \BookRecordIndex{T0372}驯服（或与 $\omega$ \BookRecordIndex{T0348}相容）的\BookRecordIndex{T0088}殆复结构
光滑族。以下一律用
\[
  \mathcal J_\tau(\Sigma;M,\omega)
  :=C^\infty\bigl(\Sigma,\mathcal J_\tau(M,\omega)\bigr)
\]
表示光滑映射空间
\[
  \Sigma\longrightarrow\mathcal J_\tau(M,\omega),
  \qquad z\longmapsto J_z.
\]
因为 $\mathcal J_\tau(M,\omega)$ 可缩，所以 $\mathcal J_\tau(\Sigma;M,\omega)$ 也可缩。
在与 $\omega$ \BookRecordIndex{T0348}相容的情形，这一结论仍成立。事实上，本节全部结果对 $\omega$-\BookRecordIndex{T0348}相容和
$\omega$-\BookRecordIndex{T0372}驯服两种情形都成立。我们用 $\omega$-\BookRecordIndex{T0372}驯服\BookRecordIndex{T0088}殆复结构来表述它们，因为这对应用更
方便。

回忆：$\mathcal J_+(M,\omega;\kappa)$ 表示 $M$ 上所有 $\kappa$-\BookRecordIndex{T0058}半正且被 $\omega$ \BookRecordIndex{T0372}驯服的
\BookRecordIndex{T0088}殆复结构之集。也就是说，$J\in\mathcal J_\tau(M,\omega)$ 属于
$\mathcal J_+(M,\omega;\kappa)$，当且仅当每个 $\omega$-能量满足 $E(v)\le\kappa$ 的
$J$-全纯球面 $v:S^2\to M$ 都有非负 Chern 数。若 $(M,\omega)$ \BookRecordIndex{T0058}半正，则由
引理~\ref{lem:6-4-8}，对每个 $\kappa>0$，集合 $\mathcal J_+(M,\omega;\kappa)$ 都非空、
开且道路连通。因此，所有满足对每个 $z\in\Sigma$ 都有
$J_z\in\mathcal J_+(M,\omega;\kappa)$ 的
$\mathbf J=\{J_z\}\in\mathcal J_\tau(\Sigma;M,\omega)$ 也构成非空开集。记
\[
  \mathcal J_+(\Sigma;M,\omega;\kappa)
  :=\Bigl\{\{J_z\}\in\mathcal J_\tau(\Sigma;M,\omega)\ \Big|\
      J_z\in\mathcal J_+(M,\omega;\kappa)\ \forall z\in\Sigma,
      \ \{J_z\}\sim\mathrm{const}\Bigr\}
\]
为其中包含常值映射的连通分支。

现在设 $\Sigma=S^2$，并在其上取标准复结构。给定
$\{J_z\}\in\mathcal J_+(S^2;M,\omega;\kappa)$、满足 $\omega(A)\le\kappa$ 的同调类
$A\in H_2(M;\mathbb Z)$、子集 $I\subset\{1,\ldots,k\}$，以及 $S^2$ 上两两不同的点所成的
元组 $\mathbf w=\{w_i\}_{i\in I}$，考察\BookRecordIndex{T0246}模空间
\[
  \mathcal M^*_{0,k}(A;\mathbf w,\{J_z\}),
\]
它由所有元组 $(u,\mathbf z)$ 组成，其中 $u:S^2\to M$ 是类 $A$ 中式~\eqref{eq:6-7-1}
的一个解，而 $\mathbf z=(z_1,\ldots,z_k)$ 是 $S^2$ 上两两不同的点所成的元组，并且当
$i\in I$ 时 $z_i=w_i$。这个\BookRecordIndex{T0246}模空间带有\BookRecordIndex{T0283}求值映射
\begin{equation}\label{eq:6-7-2}
  \operatorname{ev}:\mathcal M^*_{0,k}(A;\mathbf w,\{J_z\})\longrightarrow M^k,
  \qquad
  \operatorname{ev}(u,\mathbf z):=(u(z_1),\ldots,u(z_k)).
\end{equation}

\begin{theorem}\label{thm:6-7-1}
设 $(M,\omega)$ 是紧\BookRecordIndex{T0058}半正\BookRecordIndex{T0358}辛流形，$A\in H_2(M;\mathbb Z)$，$k\ge0$ 是整数，且
$I\subset\{1,\ldots,k\}$。设 $\mathbf w=\{w_i\}_{i\in I}$ 是 $S^2$ 上两两不同的点所成的
元组，并设 $\kappa\ge\omega(A)$。则
$\mathcal J_\tau(S^2;M,\omega)$ 中存在剩余子集
$\mathcal J_{\mathrm{reg}}(S^2;M,\omega;\mathbf w)$，使得对每个
\[
  \mathbf J=\{J_z\}_{z\in S^2}
  \in\mathcal J_{\mathrm{reg}}(S^2;M,\omega;\mathbf w)
    \cap\mathcal J_+(S^2;M,\omega;\kappa)
\]
都有以下结论。
\begin{enumerate}[label=\textup{(\roman*)},leftmargin=3em]
\item \BookRecordIndex{T0283}求值映射~\eqref{eq:6-7-2} 是维数为
\[
  \dim\mathcal M^*_{0,k}(A;\mathbf w,\{J_z\})
  =2n+2c_1(A)+2(k-\#I)
\]
的\BookRecordIndex{T0337}伪循环。
\sourcefootnote{原书把末项印作 $2(k-2\#I)$；但解空间指标为
$2n+2c_1(A)$，加入 $k$ 个可变点贡献 $2k$ 维，而固定每个定义域点只减 $2$ 维。
定理~\ref{thm:6-7-11}\textup{(i)} 在 $e(T)=0$ 时以及本节稍后的同伦\BookRecordIndex{T0246}模空间维数也都给出
$2n+2c_1(A)+2k-2\#I$，故据此补正。}
\item (i) 中\BookRecordIndex{T0337}伪循环的配边类与 $\{J_z\}$ 和 $\mathbf w$ 无关。
\item 若 $A$ 满足定理~\ref{thm:6-6-1} 的假设且 $\#I=3$，则 (i) 中的\BookRecordIndex{T0337}伪循环与
\BookRecordIndex{T0145}格罗莫夫--威滕伪循环
$\operatorname{ev}:\mathcal M^*_{0,k}(A;J)\to M^k$ 配边。
\end{enumerate}
\end{theorem}

$\mathcal J_{\mathrm{reg}}(S^2;M,\omega;\mathbf w)$ 中元素所满足的精确条件将在下面的
定义~\ref{def:6-7-10} 中列出。等我们发展出描述证明所需的适当语言后，会在本节末尾给出
证明。

上一定理只涉及 $M$ 中的曲线；只要把前几章的全部论证推广到 $J$ 依赖于 $z$ 的情形，
原则上就可以证明它。不过，若不把\BookRecordIndex{T0246}模空间 $\mathcal M(A,\Sigma;\{J_z\})$ 中式
\eqref{eq:6-7-1} 的类 $A$ 解对应到 $M$ 中的曲线，而是对应到乘积流形
\[
  \widetilde M:=\Sigma\times M
\]
中的图像，这项工作会容易得多。具体地说，可把式~\eqref{eq:6-7-1} 的解看成
$\widetilde M$ 中的\BookRecordIndex{T0334}伪全纯曲线。每个族
$\mathbf J=\{J_z\}\in\mathcal J_\tau(\Sigma;M,\omega)$ 都通过公式
\begin{equation}\label{eq:6-7-3}
  \widetilde J(z,x):=j(z)\oplus J_z(x)
  \in\operatorname{Aut}(T_z\Sigma\oplus T_xM)
\end{equation}
在 $\widetilde M$ 上给出唯一的\BookRecordIndex{T0088}殆复结构 $\widetilde J$。显然，光滑函数
$u:\Sigma\to M$ 是式~\eqref{eq:6-7-1} 的解，当且仅当其\BookRecordIndex{T0325}图像映射
$\widetilde u:\Sigma\to\widetilde M$，
\[
  \widetilde u(z):=(z,u(z)),
\]
是一条 $\widetilde J$-全纯曲线。这说明，只要 $\Sigma$ 上的复结构固定，第~4、5 章的\BookRecordIndex{T0192}紧性
论证在亏格零情形可以原封不动地搬过来，在高亏格情形也只须作少量修改。然而，第~3 章的
\BookRecordIndex{T0166}横截性论证必须因当前情形作出调整，原因有二。第一，我们只考虑 $\widetilde J$-全纯截面，
因而自然应在\BookRecordIndex{T0346}线性化算子作用于竖直向量场的空间时考察它。第二，我们只考虑与\BookRecordIndex{T0279}纤维化
$\pi:\widetilde M\to\Sigma$ \BookRecordIndex{T0348}相容的\BookRecordIndex{T0088}殆复结构 $\widetilde J$，这里“\BookRecordIndex{T0348}相容”是指 $\pi$ 全纯，且
$\pi$ 的纤维是 $\widetilde M$ 的殆复子流形。式~\eqref{eq:6-7-3} 形状的\BookRecordIndex{T0088}殆复结构满足这些
条件。

\begin{remark}\label{rem:6-7-2}
给定同调类 $A\in H_2(M;\mathbb Z)$，定义 $\widetilde A\in H_2(\widetilde M;\mathbb Z)$ 为
\[
  \widetilde A:=\iota_*A+[\Sigma\times\mathrm{pt}],
\]
其中 $\iota:M\to\widetilde M$ 表示纤维的包含映射。比较\BookRecordIndex{T0246}模空间
$\mathcal M(A,\Sigma;\{J_z\})$ 与\BookRecordIndex{T0246}模空间
$\mathcal M(\widetilde A,\Sigma;\widetilde J)/G_\Sigma$ 很有启发性；后者由类
$\widetilde A$ 的所有 $\widetilde J$-全纯代表对 $(\Sigma,j_\Sigma)$ 的复自同构群
$G_\Sigma$ 取商而得。若
$\widetilde u\in\mathcal M(\widetilde A,\Sigma;\widetilde J)$，则复合
$\phi:=\pi\circ\widetilde u$ 是 $\Sigma$ 的复自同构，因而
$\widetilde u\circ\phi^{-1}$ 是 $\pi:\widetilde M\to\Sigma$ 的
$\widetilde J$-全纯截面。换言之，存在
$u\in\mathcal M(A,\Sigma;\{J_z\})$ 和唯一的自同构 $\phi\in G_\Sigma$，使得
\[
  \widetilde u(z):=(\phi(z),u\circ\phi(z)).
\]
因此，从 $\mathcal M(A,\Sigma;\{J_z\})$ 到商空间
$\mathcal M(\widetilde A,\Sigma;\widetilde J)/G_\Sigma$ 有一个自然双射。当 $g=0$ 时，
这个商就是空间 $\mathcal M_{0,0}(\widetilde A;\widetilde J)$。在 $g>0$ 的情形，后一记号
通常留给 $\Sigma$ 的复结构可在 Teichmüller 空间中自由变化的空间，而我们的方法只处理
$j_\Sigma$ 固定的情形。Teichmüller 空间中的这种变化，正是下面注~\ref{rem:6-7-5} 所提到
的算子 $D_{\widetilde u}$ 的余核所具有的几何意义。还要注意，若
$(\Sigma,j_\Sigma)$ 是亏格 $g>1$ 的一般黎曼曲面，则它没有自同构
\sourcefootnote{若这里的“没有自同构”意指没有非平凡自同构，则原书的范围 $g>1$ 并不
包含亏格 $2$：每个亏格 $2$ 曲面都有超椭圆对合；相应的通有无非平凡自同构结论应取
$g>2$。若这里只意指没有连续的或无穷小的自同构，则原范围成立。}；而对一般的 $j$，
$(\mathbb T^2,j)$ 的自同构群有两个连通分支。
\end{remark}

\paragraph{图像\BookRecordIndex{T0246}模空间。}
给定紧黎曼曲面 $(\Sigma,j_\Sigma)$、\BookRecordIndex{T0088}殆复结构
$\{J_z\}\in\mathcal J_\tau(\Sigma;M,\omega)$ 和同调类 $A\in H_2(M;\mathbb Z)$，记表示
类 $A$ 的 $\{J_z\}$-全纯曲线 $u:\Sigma\to M$ 所成的\BookRecordIndex{T0246}模空间为
\[
  \mathcal M(A,\Sigma;\{J_z\})
  :=\bigl\{u:\Sigma\to M\mid u\text{ 满足式~\eqref{eq:6-7-1}},\ [u]=A\bigr\}.
\]
当 $\Sigma=S^2$ 时，简记
\[
  \mathcal M(A;\{J_z\}):=\mathcal M(A,S^2;\{J_z\}).
\]
要理解这个\BookRecordIndex{T0246}模空间，必须考察\BookRecordIndex{T0346}线性化算子
\[
  D_u=D_{u,\mathbf J}:\Omega^0(\Sigma,u^*TM)
  \longrightarrow\Omega^{0,1}(\Sigma,u^*TM).
\]
如第~\ref{sec:3-1} 节所述，该算子由公式
\begin{equation}\label{eq:6-7-4}
  D_u\xi
  :=\frac12\bigl(\nabla\xi+J_z(u)\nabla\xi\circ j_\Sigma\bigr)
    -\frac12J_z(u)(\nabla_\xi J_z)\partial_J(u)
\end{equation}
给出；但现在 $\nabla=\nabla^z$ 表示由 $\omega$ 和 $J_z$ 依式~\eqref{eq:2-1-1} 确定的
黎曼度量
\[
  \langle\cdot,\cdot\rangle_z
  :=\frac12\bigl(\omega(\cdot,J_z\cdot)-\omega(J_z\cdot,\cdot)\bigr)
\]
的 \BookRecordIndex{T0034}Levi--Civita 联络。因此，对每个 $z\in\Sigma$，我们都用联络 $\nabla^z$ 和\BookRecordIndex{T0088}殆复结构
$J_z$ 来计算 $(D_u\xi)(z)$。同以前一样，这是实线性\BookRecordIndex{T0210}柯西--黎曼算子，所以由定理 C.1.10，
它是\BookRecordIndex{T0128}弗雷德霍姆算子，且
\[
  \operatorname{index}D_u=n(2-2g)+2c_1(A),
\]
其中 $g$ 是 $\Sigma$ 的亏格。

\begin{remark}\label{rem:6-7-3}
引入这个算子的动机与第~\ref{sec:3-1} 节相同。对每个
$u\in C^\infty(\Sigma,M)$，都有映射
\[
  \mathcal F_u:\Omega^0(\Sigma,u^*TM)\longrightarrow\Omega^{0,1}(\Sigma,u^*TM),
\]
定义为
\[
  \mathcal F_u(\xi):=\Phi_u(\xi)^{-1}
  \bar\partial_{\mathbf J}\bigl(\exp_u(\xi)\bigr),
\]
其中 $\exp_{u(z)}:T_{u(z)}M\to M$ 表示度量 $\langle\cdot,\cdot\rangle_z$ 的指数映射，而
\[
  \Phi_{u(z)}(\xi(z)):T_{u(z)}M\longrightarrow T_{\exp_{u(z)}(\xi(z))}M
\]
表示关于\BookRecordIndex{T0052}埃尔米特联络 $\widetilde\nabla^z$ 的\BookRecordIndex{T0263}平行移动；如第~\ref{sec:2-1} 节所述，这一联络
由 $\nabla^z$ 得到。在这些约定下，
\[
  D_u=d\mathcal F_u(0).
\]
证明与命题~\ref{prop:3-1-1} 的证明逐字相同，故从略。与第~\ref{sec:3-1} 节一样，这套设置
可推广到适当的索伯列夫完备化。
\end{remark}

\begin{remark}\label{rem:6-7-4}
若每个 $J_z$ 都与 $\omega$ \BookRecordIndex{T0348}相容，则 $D_u$ 的复线性部分把每个
$\xi\in\Omega^0(\Sigma,u^*TM)$ 映到 $1$-形式
$\widetilde\nabla\xi\in\Omega^1(\Sigma,u^*TM)$ 的复线性部分，其中
\[
  \widetilde\nabla\xi
  :=\nabla\xi-\frac12J(\nabla_{du}J)\xi-\frac12J\dot J\xi.
\]
这里的线性映射
\[
  T\Sigma\longrightarrow\Omega^0(M,\operatorname{End}(TM)),
  \qquad \zeta\longmapsto\dot J_\zeta,
\]
表示映射 $z\mapsto J_z$ 的微分。于是，把 $D_u$ 分解成复线性部分和复反线性部分
（如注~\ref{rem:3-1-3} 中的式~\eqref{eq:3-1-6}），得到
\[
  D_u\xi=(\widetilde\nabla\xi)^{0,1}
    +\frac14N_J(\xi,\partial_J(u))
    +\frac12(J\dot J\xi)^{0,1}.
\]
这里，对每个 $z\in\Sigma$，$N_J$ 表示 $J_z$ 的 \BookRecordIndex{T0042}Nijenhuis 张量。这个公式的证明只须直接
推广引理 C.7.3 证明中的论证，故从略。
\end{remark}

\begin{remark}\label{rem:6-7-5}
设 $u:\Sigma\to M$ 是任意光滑映射，并以
$\widetilde u:\Sigma\to\widetilde M$ 表示由
$\widetilde u(z):=(z,u(z))$ 定义的映射。与第~\ref{sec:3-1} 节一样，映射
$\widetilde u(z)=(z,u(z))$ 对应一个算子
\[
  D_{\widetilde u}:\Omega^0(\Sigma,\widetilde u^*T\widetilde M)
  \longrightarrow\Omega^{0,1}(\Sigma,\widetilde u^*T\widetilde M).
\]
其指标为
\[
  \operatorname{index}D_{\widetilde u}
  =(n+1)(2-2g)+2c_1(\widetilde A)
  =\operatorname{index}D_u+6-6g.
\]
它与 $D_u$ 的关系如下。复向量丛之间有自然同构
\[
  \widetilde u^*T\widetilde M\cong T\Sigma\oplus u^*TM.
\]
一般而言，\BookRecordIndex{T0308}竖直切丛和水平切丛都不在 $D_{\widetilde u}$ 下不变。不过，若 $u$ 是式
\eqref{eq:6-7-1} 的解，且 $\widetilde J$ 由式~\eqref{eq:6-7-3} 给出，则在局部坐标中考察
算子可得
\begin{equation}\label{eq:6-7-5}
  D_{\widetilde u}(\zeta,\xi)
  =\left(\bar\partial\zeta,
    D_u\xi-\frac12J_z(u)\dot J_\zeta(u)\partial_J(u)\right)
\end{equation}
其中 $\zeta\in\Omega^0(\Sigma,T\Sigma)$，而
$\xi\in\Omega^0(\Sigma,u^*TM)$。式~\eqref{eq:6-7-5} 表明，竖直子丛
$u^*TM\subset\widetilde u^*T\widetilde M$ 在 $D_{\widetilde u}$ 下不变，且 $D_u$ 是
$D_{\widetilde u}$ 在该子丛上的限制。此外，若 $\Sigma=S^2$，而
$u:S^2\to M$ 是式~\eqref{eq:6-7-1} 的解，则 $D_u$ 满射当且仅当
$D_{\widetilde u}$ 满射。要说明这一点，注意算子
\[
  \bar\partial:\Omega^0(S^2,TS^2)\longrightarrow\Omega^{0,1}(S^2,TS^2)
\]
满射，其核由全纯向量场组成；而只要 $u:S^2\to M$ 是式~\eqref{eq:6-7-1} 的解，且
$\zeta\in\Vect(S^2)$ 是全纯向量场，就有
\[
  \frac12J_z(u)\dot J_\zeta(u)\partial_J(u)=D_u(du\circ\zeta).
\]
所以，当 $\bar\partial\zeta=0$ 时，项
$J_z(u)\dot J_\zeta(u)\partial_J(u)$ 对 $D_{\widetilde u}$ 的像没有贡献任何
$D_u$ 的像中尚未包含的东西。另一方面，当 $g>0$ 时，$\bar\partial$ 有非零余核，故
$D_{\widetilde u}$ 从不满射。在亏格零情形，$D_{\widetilde u}$ 的额外高维核恰好给出
$\widetilde u$ 在\BookRecordIndex{T0413}重参数化群 $G=\operatorname{PSL}(2,\mathbb C)$ 作用下轨道的切空间。
\end{remark}

一个重要观察是：第~3 章的\BookRecordIndex{T0166}横截性论证对式~\eqref{eq:6-7-1} 的所有解 $u$ 都成立，而无论
它们是否某处单射。直观地说，这是因为图像总是嵌入的。作为热身，我们证明在
$\Sigma=S^2$ 时，式~\eqref{eq:6-7-1} 的常值解都是正则的。

\begin{lemma}\label{lem:6-7-6}
设 $(\Sigma,j_\Sigma)$ 是闭连通黎曼曲面，$u:\Sigma\to M$ 是取值为 $x\in M$ 的常值映射，
而 $\mathbf J=\{J_z\}$ 是 $M$ 上被 $\omega$ \BookRecordIndex{T0372}驯服的\BookRecordIndex{T0088}殆复结构光滑族。令 $D_{x,\mathbf J}$
为式~\eqref{eq:6-7-4} 定义的算子。则 $D_{x,\mathbf J}$ 的核由常值映射
$\xi:\Sigma\to T_xM$ 组成。此外，$D_{x,\mathbf J}$ 满射当且仅当 $\Sigma$ 的亏格为零。
\end{lemma}

\begin{proof}
算子
\[
  D_{x,\mathbf J}:\Omega^0(\Sigma,T_xM)\longrightarrow\Omega^{0,1}(\Sigma,T_xM)
\]
由
\[
  (D_{x,\mathbf J}\xi)(z)
  =\frac12\bigl(\nabla\xi(z)+J(z)\nabla\xi(z)\circ j_\Sigma(z)\bigr)
\]
给出，其中 $\xi:\Sigma\to T_xM$，$z\in\Sigma$。若 $J$ 与 $z$ 无关，第一个断言显然成立，
因为它等价于说 $\Sigma$ 上任意全纯函数都是常值。在一般情形，论证如下。由于
$u(z)\equiv x$，$\xi$ 的协变导数就是普通导数。以 $|\nabla\xi(z)|$ 表示线性映射
$\nabla\xi(z):T_z\Sigma\to T_xM$ 的范数，其中 $T_xM$ 上的内积由 $\omega$ 和 $J_z$ 依
式~\eqref{eq:2-1-1} 确定（见习题 2.2.3）。于是
\begin{equation}\label{eq:6-7-6}
  D_x\xi=0
  \quad\Longrightarrow\quad
  \frac12|\nabla\xi|^2\,\mathrm{dvol}_\Sigma=\xi^*\omega_x.
\end{equation}
\sourcefootnote{原书在式~\eqref{eq:6-7-6} 中把体积形式的下标印作 $S^2$，但本引理的定义域
是任意闭黎曼曲面 $\Sigma$，紧随其后的局部计算也使用 $\mathrm{dvol}_\Sigma$。原书随后又把
共形度量 $\lambda^2(ds^2+dt^2)$ 误称为体积形式；正文分别改为正确的
$\mathrm{dvol}_\Sigma$，并明确相应体积形式为 $\lambda^2ds\wedge dt$。}
为验证这一点，取 $\Sigma$ 上的\BookRecordIndex{T0149}共形坐标 $s+it$，并设这些坐标下的共形度量为
$\lambda^2(ds^2+dt^2)$，因而
$\mathrm{dvol}_\Sigma=\lambda^2ds\wedge dt$。若 $D_x\xi=0$，则
$\partial_s\xi+J_{s,t}(x)\partial_t\xi=0$，从而
\begin{align*}
  \lambda^2|\nabla\xi|^2
  &=|\partial_s\xi|^2+|\partial_t\xi|^2\\
  &=\omega(\partial_s\xi,J_{s,t}(x)\partial_s\xi)
    +\omega(\partial_t\xi,J_{s,t}(x)\partial_t\xi)\\
  &=2\omega(\partial_s\xi,\partial_t\xi).
\end{align*}
这就证明了式~\eqref{eq:6-7-6}。在 $\Sigma$ 上积分可知，每个
$\xi\in\ker D_{x,\mathbf J}$ 都满足 $\nabla\xi\equiv0$，因而是常值。这证明了第一个断言。
对第二个断言，由定理 C.1.10，算子 $D_{x,\mathbf J}$ 的\BookRecordIndex{T0129}弗雷德霍姆指标为
$n(2-2g)$，其中 $g$ 是 $\Sigma$ 的亏格。刚才已经证明它的核维数为 $2n$。所以
$D_{x,\mathbf J}$ 满射当且仅当 $g=0$。
\end{proof}

相应的\BookRecordIndex{T0088}殆复结构\emph{正则族}之集记作\footnote{这里的
$\mathcal J_{\mathrm{reg}}(\Sigma;A)$ 不等于第~\ref{sec:3-1} 节所考察的集合
$\mathcal J_{\mathrm{reg}}(A,\Sigma)$；后者由所有与 $z$ 无关、且使类 $A$ 中的单纯
\BookRecordIndex{T0003}$J$-全纯曲线 $u:\Sigma\to M$ 正则的\BookRecordIndex{T0088}殆复结构 $J$ 组成。本节的全部讨论都针对依赖于
$z$ 的\BookRecordIndex{T0088}殆复结构，二者记号相近不应造成混淆。}
\[
  \mathcal J_{\mathrm{reg}}(\Sigma;A)
  :=\Bigl\{\{J_z\}\in\mathcal J_\tau(\Sigma;M,\omega)\ \Big|\
      u\in\mathcal M(A,\Sigma;\{J_z\})\Longrightarrow D_u\text{ 满射}\Bigr\}.
\]
给定 $\mathcal J_\tau(\Sigma;M,\omega)$ 中的两个\BookRecordIndex{T0088}殆复结构族
$\mathbf J^0=\{J_z^0\}$ 与 $\mathbf J^1=\{J_z^1\}$，以
$\mathcal J_\tau(\Sigma;M,\omega;\mathbf J^0,\mathbf J^1)$ 表示从 $\mathbf J^0$ 到
$\mathbf J^1$ 的光滑同伦
\[
  [0,1]\times\Sigma\longrightarrow\mathcal J_\tau(M,\omega),
  \qquad (\lambda,z)\longmapsto J_z^\lambda
\]
所成的空间。给定这样的同伦，考察\BookRecordIndex{T0246}模空间
\[
  \mathcal W(A,\Sigma;\{J_z^\lambda\})
  :=\bigl\{(\lambda,u)\mid 0\le\lambda\le1,
      \ u\in\mathcal M(A,\Sigma;\mathbf J^\lambda)\bigr\}.
\]
若 $\{J_z^0\},\{J_z^1\}\in\mathcal J_{\mathrm{reg}}(\Sigma;A)$，并且对每个
$(\lambda,u)\in\mathcal W(A,\Sigma;\{J_z^\lambda\})$ 都有
\[
  \Omega^{0,1}(\Sigma,u^*TM)
  =\operatorname{im}D_u
    +\mathbb R(\partial_\lambda J_z^\lambda(u))du\circ j_\Sigma,
\]
就称同伦 $\{J_z^\lambda\}$ \emph{（关于 $A$）正则}。从 $\mathbf J^0$ 到 $\mathbf J^1$
的\BookRecordIndex{T0407}正则同伦之集记作
$\mathcal J_{\mathrm{reg}}(\Sigma;A;\mathbf J^0,\mathbf J^1)$。

\begin{proposition}\label{prop:6-7-7}
设 $A\ne0$ 或 $\Sigma=S^2$。
\begin{enumerate}[label=\textup{(\roman*)},leftmargin=3em]
\item $\mathcal J_{\mathrm{reg}}(\Sigma;A)$ 是
$\mathcal J_\tau(\Sigma;M,\omega)$ 中的\BookRecordIndex{T0300}剩余集。
\item 若 $\{J_z\}\in\mathcal J_{\mathrm{reg}}(\Sigma;A)$，则
$\mathcal M(A,\Sigma;\{J_z\})$ 是光滑定向流形，其维数为
\[
  \dim\mathcal M(A,\Sigma;\{J_z\})=n(2-2g)+2c_1(A).
\]
\item 设
$\{J_z^0\},\{J_z^1\}\in\mathcal J_{\mathrm{reg}}(\Sigma;A)$。则
$\mathcal J_{\mathrm{reg}}(\Sigma;A;\mathbf J^0,\mathbf J^1)$ 是
$\mathcal J_\tau(\Sigma;M,\omega;\mathbf J^0,\mathbf J^1)$ 中的\BookRecordIndex{T0300}剩余集。此外，若
$\{J_z^\lambda\}\in
\mathcal J_{\mathrm{reg}}(\Sigma;A;\mathbf J^0,\mathbf J^1)$，则
$\mathcal W(A,\Sigma;\{J_z^\lambda\})$ 是光滑定向带边流形，且
\[
  \partial\mathcal W(A,\Sigma;\{J_z^\lambda\})
  =\mathcal M(A,\Sigma;\{J_z^1\})
    \cup\bigl(-\mathcal M(A,\Sigma;\{J_z^0\})\bigr).
\]
\end{enumerate}
\end{proposition}

\begin{proof}
当 $\Sigma=S^2$ 且 $A=0$ 时，结论显然来自引理~\ref{lem:6-7-6}。当 $A\ne0$ 时，证明遵循
定理~\ref{thm:3-1-6} 和定理~\ref{thm:3-1-8} 的同一模式。先证明\BookRecordIndex{T0268}普适模空间是流形
（对应于命题~\ref{prop:3-2-1}），再考察从这个普适空间到\BookRecordIndex{T0088}殆复结构参数空间
$\mathcal J_\tau(\Sigma;M,\omega)$ 的投影，以证明 (i) 和 (ii)。

更准确地说，记 $\mathcal J^\ell=\mathcal J_\tau^\ell(\Sigma;M,\omega)$ 为 $M$ 上被
$\omega$ \BookRecordIndex{T0372}驯服的 $C^\ell$ \BookRecordIndex{T0088}殆复结构族所成的空间；这意味着 $\widetilde M$ 上的诱导复结构
$\widetilde J$ 属于 $C^\ell$。考察\BookRecordIndex{T0268}普适模空间
\[
  \mathcal M(A,\Sigma;\mathcal J^\ell)
  :=\bigl\{(u,\{J_z\})\mid \{J_z\}\in\mathcal J^\ell,
      \ u\in\mathcal M(A,\Sigma;\{J_z\})\bigr\}.
\]
我们证明，该空间是
$W^{1,p}(\Sigma,M)\times\mathcal J^\ell$ 的 $C^{\ell-1}$ 巴拿赫子流形。依照命题
~\ref{prop:3-2-1} 证明中的论证，须证明\BookRecordIndex{T0346}线性化算子
\[
  W^{1,p}(\Sigma,u^*TM)\times T_{\{J_z\}}\mathcal J^\ell
  \longrightarrow\Omega^{0,1}(\Sigma,u^*TM),
  \qquad
  (\xi,\{Y_z\})\longmapsto
  D_u\xi+\frac12Y_z(u)du\circ j_\Sigma
\]
在每个 $(u,\{J_z\})\in\mathcal M(A,\Sigma;\mathcal J^\ell)$ 处都满射。为此，取满足
$1/p+1/q=1$ 的 $q>1$，并设
\[
  \eta\in L^q\bigl(\Sigma,T^*\Sigma\otimes_{\{J_z\}}u^*TM\bigr)
\]
零化这个算子的像。于是 $\eta\in W^{\ell,q}$
\sourcefootnote{原书把\BookRecordIndex{T0329}椭圆正则性提升后的空间印作 $W^{\ell,p}$；但 $\eta$ 起初属于
$L^q$，其中 $q$ 是 $p$ 的对偶指数，作用于伴随算子 $D_u^*$ 的\BookRecordIndex{T0329}椭圆正则性保持这一积分
指数，故应为 $W^{\ell,q}$。}，且
\[
  D_u^*\eta=0,
\]
并且对每个 $\{Y_z\}\in T_{\{J_z\}}\mathcal J^\ell$ 都有
\begin{equation}\label{eq:6-7-7}
  \int_\Sigma\bigl\langle\eta,Y_z(u)du\circ j_\Sigma\bigr\rangle
  \,\mathrm{dvol}_\Sigma=0.
\end{equation}
由于 $A\ne0$，函数 $u$ 不可能是常值。引理~\ref{lem:2-4-1} 的\BookRecordIndex{T0233}临界点分析可以照搬过来，
并说明 $du$ 只可能在有限多个点处消失。我们证明：对每个满足 $du(z)\ne0$ 的
$z\in\Sigma$，都有 $\eta(z)=0$。反设不然，即对某个 $z_0\in\Sigma$，有
$du(z_0)\ne0$ 且 $\eta(z_0)\ne0$。于是，由引理~\ref{lem:3-2-2}，存在\BookRecordIndex{T0343}无穷小殆复结构
$\{Y_z\}\in T_{\{J_z\}}\mathcal J^\ell$，使函数
\[
  z\longmapsto\bigl\langle\eta(z),Y_z(u(z))du(z)\circ j_\Sigma\bigr\rangle
\]
在 $z_0$ 处为正，因而在 $z_0$ 的某个邻域 $U\subset\Sigma$ 中为正。取支集位于 $U$ 内且
$\beta(z_0)=1$ 的光滑\BookRecordIndex{T0188}截断函数 $\beta:\Sigma\to[0,1]$。则\BookRecordIndex{T0343}无穷小殆复结构
$z\mapsto\beta(z)Y_z$ 违反条件~\eqref{eq:6-7-7}。这个矛盾说明 $\eta$ 几乎处处为零，从而
$\eta\equiv0$。因此，\BookRecordIndex{T0346}线性化算子满射，\BookRecordIndex{T0268}普适模空间
$\mathcal M(A,\Sigma;\mathcal J^\ell)$ 是\BookRecordIndex{T0213}可分的 $C^{\ell-1}$ \BookRecordIndex{T0053}巴拿赫流形。证明的其余部分
与定理~\ref{thm:3-1-6} 和定理~\ref{thm:3-1-8} 的证明逐字相同，故从略。
\end{proof}

\paragraph{纤维中的 $J$-全纯球面。}
\BookRecordIndex{T0246}模空间 $\mathcal M(A,\Sigma;\{J_z\})$ 的 Gromov--Kontsevich 紧化会涉及 $M$ 中的
$J_z$-全纯球面。从\BookRecordIndex{T0141}格罗莫夫收敛定义里的重标度性质可以清楚看出这一点。设
$u_\nu:\Sigma\to M$ 是式~\eqref{eq:6-7-1} 的一列解，其一阶导数在
$z\in\Sigma$ 附近发散；又设 $\phi_\nu:B_{R_\nu}\to\Sigma$ 是一列全纯嵌入，它在每个
紧子集上收敛到 $z$。那么，序列 $u_\nu\circ\phi_\nu$ 的极限（若存在）将是一条
$J_z$-全纯球面。因此，我们必须考察所有二元组 $(z,v)$ 所成的\BookRecordIndex{T0246}模空间，其中
$v:S^2\to M$ 是 $J_z$-全纯球面。也可以把它们看成
$\widetilde M=\Sigma\times M$ 中竖直的 $\widetilde J$-全纯球面。

给定球面同调类 $A\in H_2(M;\mathbb Z)$ 和\BookRecordIndex{T0088}殆复结构
$\mathbf J=\{J_z\}\in\mathcal J_\tau(\Sigma;M,\omega)$，考察 $\widetilde M$ 中表示类 $A$
的单纯竖直 $J_z$-全纯球面\BookRecordIndex{T0246}模空间
\[
  \mathcal M^{\mathrm{Vert},*}(A;\{J_z\})
  :=\bigl\{(z,v)\mid z\in\Sigma,\ v\in\mathcal M^*(A;J_z)\bigr\}.
\]
相应的\BookRecordIndex{T0346}线性化算子
\begin{equation}\label{eq:6-7-8}
  D_{z,v}:T_z\Sigma\oplus\Omega^0(S^2,v^*TM)
  \longrightarrow\Omega^{0,1}(S^2,v^*TM)
\end{equation}
由
\[
  D_{z,v}(\zeta,\xi)
  =D_{v,J_z}\xi-\frac12J_z(v)\dot J_\zeta(v)\partial_{J_z}(v)
\]
给出，其中线性映射
\[
  T\Sigma\longrightarrow\Omega^0(M,\operatorname{End}(TM)),
  \qquad \zeta\longmapsto\dot J_\zeta
\]
表示映射 $z\mapsto J_z$ 的微分，如注~\ref{rem:6-7-4} 所述。为说明这个公式，先固定两条
道路
\[
  \mathbb R\longrightarrow\Sigma:\lambda\longmapsto z(\lambda),
  \qquad
  \mathbb R\longrightarrow M:\lambda\longmapsto x(\lambda).
\]
道路 $\lambda\mapsto z(\lambda)$ 在拉回切丛 $x^*TM$ 上定义一个复结构；它在
$T_{x(\lambda)}M$ 上由 $J_{z(\lambda)}(x(\lambda))$ 给出。对沿 $x$ 的向量场
$\xi(\lambda)\in T_{x(\lambda)}M$，公式
\[
  \widetilde\nabla_\lambda\xi
  :=\nabla_\lambda\xi
    -\frac12J(\nabla_{\dot x(\lambda)}J)\xi
    -\frac12J\dot J_{\dot z(\lambda)}\xi
\]
在 $x^*TM$ 上定义一个保持此复结构的联络。这里，$\nabla_\lambda\xi$ 表示关于由
$J_{z(\lambda)}$ 关联的 \BookRecordIndex{T0034}Levi--Civita 联络，沿 $\lambda$ 方向对 $\xi$ 所取的协变导数。
给定 $\zeta\in T_z\Sigma$ 和 $\xi\in T_xM$，以
\[
  \Phi_{z,x}(\zeta,\xi):T_xM\longrightarrow T_{\exp_x(\xi)}M
\]
表示沿道路 $\lambda\mapsto\exp_x(\lambda\xi)$ 的\BookRecordIndex{T0263}平行移动；所用联络是与道路
$\lambda\mapsto\exp_z(\lambda\zeta)$ 对应的 $\widetilde\nabla$。于是，与注
~\ref{rem:6-7-3} 一样，有映射
\[
  \mathcal F_{z,v}:T_z\Sigma\oplus\Omega^0(S^2,v^*TM)
  \longrightarrow\Omega^{0,1}(S^2,v^*TM),
\]
它把 $(\zeta,\xi)$ 映为 $(0,1)$-形式
\[
  \mathcal F_{z,v}(\zeta,\xi)
  :=\Phi_{z,v}(\zeta,\xi)^{-1}
    \bar\partial_{J_{\exp_z(\zeta)}}\bigl(\exp_v(\xi)\bigr).
\]
用这个记号，可以像命题~\ref{prop:3-1-1} 那样证明
$D_{z,v}=d\mathcal F_{z,v}(0)$。

\begin{remark}\label{rem:6-7-8}
给定光滑映射 $v:S^2\to M$ 和点 $z\in\Sigma$，定义
\[
  \widetilde v:S^2\longrightarrow\widetilde M,
  \qquad \widetilde v(z'):=(z,v(z')).
\]
则 $v$ 是 $M$ 中的 $J_z$-全纯球面，当且仅当 $\widetilde v$ 是 $\widetilde M$ 中的
$\widetilde J$-全纯球面。相应地，有\BookRecordIndex{T0346}线性化算子
\[
  D_{\widetilde v}:\Omega^0(S^2,\widetilde v^*T\widetilde M)
  \longrightarrow\Omega^{0,1}(S^2,\widetilde v^*T\widetilde M).
\]
$D_{\widetilde v}$ 和 $D_{z,v}$ 都是\BookRecordIndex{T0128}弗雷德霍姆算子，且指标相同：
\[
  \operatorname{index}D_{z,v}
  =\operatorname{index}D_{\widetilde v}
  =2n+2+2c_1(A).
\]
此外，两者有如下关系。以 $\pi:\widetilde M\to\Sigma$ 表示自然投影，则
\[
  d\pi(\widetilde v)\circ D_{\widetilde v}
  =\bar\partial\circ d\pi(\widetilde v),
\]
其中
\[
  \bar\partial:\Omega^0(S^2,T_z\Sigma)
  \longrightarrow\Omega^{0,1}(S^2,T_z\Sigma)
\]
是平凡丛 $S^2\times T_z\Sigma\to S^2$ 上的\BookRecordIndex{T0210}柯西--黎曼算子。这个恒等式说明，算子
$D_{\widetilde v}$ 把子空间
\[
  \bigl\{\widetilde\xi\in\Omega^0(S^2,\widetilde v^*T\widetilde M)
    \mid d\pi(\widetilde v)\widetilde\xi=\mathrm{const}\bigr\}
\]
映到 $d\pi(\widetilde v)$ 的核中，而 $D_{\widetilde v}$ 在此子空间上的限制恰是
$D_{z,v}$。由于 $\bar\partial$ 满射，且其核是常值函数
$S^2\to T_z\Sigma$ 所成的空间，$D_{\widetilde v}$ 与 $D_{z,v}$ 有相同的核以及同构的
余核。特别地，$D_{\widetilde v}$ 满射当且仅当 $D_{z,v}$ 满射。因此，竖直球面
$\widetilde v$ 作为乘积 $\widetilde M$ 中的曲线正则，当且仅当它在 Buse
\srcite{54} 所定义的意义下\BookRecordIndex{T0071}参数正则。另一个相关例子见例 8.4.3。
\end{remark}

下面的命题明确给出竖直 $J_z$-全纯球面的相应\BookRecordIndex{T0166}横截性结果。用
\[
  \mathcal J_{\mathrm{reg}}^{\mathrm{Vert}}(\Sigma;A)
  :=\Bigl\{\{J_z\}\in\mathcal J_\tau(\Sigma;M,\omega)\ \Big|\
      z\in\Sigma,\ v\in\mathcal M^*(A;J_z)
      \Longrightarrow D_{z,v}\text{ 满射}\Bigr\}
\]
表示类 $A$ 中竖直 $J_z$-全纯球面的正则\BookRecordIndex{T0088}殆复结构族之空间。同伦
$\{J_z^\lambda\}\in
\mathcal J_\tau(\Sigma;M,\omega;\mathbf J^0,\mathbf J^1)$ 确定\BookRecordIndex{T0246}模空间
\[
  \mathcal W^{\mathrm{Vert},*}(A;\{J_z^\lambda\})
  :=\bigl\{(\lambda,z,v)\mid 0\le\lambda\le1,\ z\in\Sigma,
      \ v\in\mathcal M^*(A;J_z^\lambda)\bigr\}.
\]
\sourcefootnote{原书在这里把参数区间印作 $0\le\lambda<1$；但同伦的定义域是闭区间
$[0,1]$，且命题~\ref{prop:6-7-9}\textup{(iii)} 随即把 $\lambda=1$ 的\BookRecordIndex{T0246}模空间列为配边的一侧，
故补上端点 $\lambda=1$。}
若
$\mathbf J^0,\mathbf J^1\in\mathcal J_{\mathrm{reg}}^{\mathrm{Vert}}(\Sigma;A)$，并且对
每个 $(\lambda,z,v)\in\mathcal W^{\mathrm{Vert},*}(A;\{J_z^\lambda\})$ 都有
\begin{equation}\label{eq:6-7-9}
  \Omega^{0,1}(S^2,v^*TM)
  =\operatorname{im}D_{z,v}
    +\mathbb R(\partial_\lambda J_z^\lambda(v))dv\circ j_{S^2},
\end{equation}
\sourcefootnote{原书把式~\eqref{eq:6-7-9} 左端的定义域印作 $\Sigma$；但
$v:S^2\to M$，式~\eqref{eq:6-7-8} 中 $D_{z,v}$ 的靶空间以及右端
$dv\circ j_{S^2}$ 都定义在 $S^2$ 上，故补正为 $S^2$。}
就称该同伦\emph{（关于 $A$）\BookRecordIndex{T0311}竖直正则}。从 $\mathbf J^0$ 到 $\mathbf J^1$ 的\BookRecordIndex{T0311}竖直正则
同伦之集记作
$\mathcal J_{\mathrm{reg}}^{\mathrm{Vert}}
(\Sigma;A;\mathbf J^0,\mathbf J^1)$。下面的结果可以看作定理~\ref{thm:3-1-6} 和
定理~\ref{thm:3-1-8} 的参数化版本。证明只是直接改写第~\ref{sec:3-2} 节中的论证，故从略。

\begin{proposition}\label{prop:6-7-9}
设 $A\in H_2(M;\mathbb Z)$ 是球面同调类。
\begin{enumerate}[label=\textup{(\roman*)},leftmargin=3em]
\item $\mathcal J_{\mathrm{reg}}^{\mathrm{Vert}}(\Sigma;A)$ 是
$\mathcal J_\tau(\Sigma;M,\omega)$ 中的\BookRecordIndex{T0300}剩余集。
\item 若 $\{J_z\}\in\mathcal J_{\mathrm{reg}}^{\mathrm{Vert}}(\Sigma;A)$，则
$\mathcal M^{\mathrm{Vert},*}(A;\{J_z\})$ 是光滑定向流形，其维数为
\[
  \dim\mathcal M^{\mathrm{Vert},*}(A;\{J_z\})=2n+2+2c_1(A).
\]
\item 设
$\mathbf J^0,\mathbf J^1\in\mathcal J_{\mathrm{reg}}^{\mathrm{Vert}}(\Sigma;A)$。则
$\mathcal J_{\mathrm{reg}}^{\mathrm{Vert}}
(\Sigma;A;\mathbf J^0,\mathbf J^1)$ 是
$\mathcal J_\tau(\Sigma;M,\omega;\mathbf J^0,\mathbf J^1)$ 中的\BookRecordIndex{T0300}剩余集。此外，若
$\{J_z^\lambda\}\in\mathcal J_{\mathrm{reg}}^{\mathrm{Vert}}
(\Sigma;A;\mathbf J^0,\mathbf J^1)$，则
$\mathcal W^{\mathrm{Vert},*}(A;\{J_z^\lambda\})$ 是光滑定向带边流形，且
\[
  \partial\mathcal W^{\mathrm{Vert},*}(A;\{J_z^\lambda\})
  =\mathcal M^{\mathrm{Vert},*}(A;\{J_z^1\})
    \cup\bigl(-\mathcal M^{\mathrm{Vert},*}(A;\{J_z^0\})\bigr).
\]
\end{enumerate}
\end{proposition}

\paragraph{\BookRecordIndex{T0340}稳定映射。}
为便于叙述，现在把注意力限制到 $\Sigma=S^2$ 的情形，这样就可以不作任何修改地直接应用
第~5 章的记号和结果。第~8 章将在 Hamilton 扰动的背景下详细讨论怎样把这些结果推广到
高亏格情形。

固定 $M$ 上被 $\omega$ \BookRecordIndex{T0372}驯服的\BookRecordIndex{T0088}殆复结构光滑族
$\mathbf J=\{J_z\}_{z\in S^2}\in\mathcal J_\tau(S^2;M,\omega)$，并以
$\widetilde J$ 表示乘积流形 $\widetilde M:=S^2\times M$ 上相应的\BookRecordIndex{T0088}殆复结构。于是，式
~\eqref{eq:6-7-1} 的解以及 $M$ 中的 $J_z$-全纯球面，都可以表成 $\widetilde M$ 中的
$\widetilde J$-全纯球面。设 $T=(T,E,\Lambda)$ 是一棵带 $k$ 个标号、具有\BookRecordIndex{T0321}特殊顶点（即根）
$0\in T$ 的树，并设 $\{A_\alpha\}_{\alpha\in T}$ 是 $H_2(M;\mathbb Z)$ 中的一族同调类。
令 $\iota:M\to\widetilde M$ 为纤维的包含映射，并记
\[
  \widetilde A_\alpha
  :=\begin{cases}
    \iota_*A_0+[S^2\times\mathrm{pt}],&\alpha=0,\\
    \iota_*A_\alpha,&\alpha\ne0.
  \end{cases}
\]
于是
\[
  \pi_*\widetilde A_\alpha
  =\begin{cases}
    [S^2],&\alpha=0,\\
    0,&\alpha\ne0,
  \end{cases}
\]
其中 $\pi:\widetilde M\to S^2$ 表示投影。对每个 $\alpha\in T\setminus\{0\}$，类
$\widetilde A_\alpha$ 中的 $\widetilde J$-全纯球面 $\widetilde u_\alpha$ 与二元组
$(z_\alpha,u_\alpha)$ 一一对应，其中 $u_\alpha$ 是类 $A_\alpha$ 中的
$J_{z_\alpha}$-全纯球面。同样，满足 $\pi\circ\widetilde u_0=\operatorname{id}$ 的类
$\widetilde A_0$ 中 $\widetilde J$-全纯球面 $\widetilde u_0$，与类 $A_0$ 中式
~\eqref{eq:6-7-1} 的解 $u_0$ 一一对应。

现在考察第~5 章定义的单纯\BookRecordIndex{T0340}稳定映射\BookRecordIndex{T0246}模空间
\[
  \mathcal M^*_{0,T}(\widetilde M,\{\widetilde A_\alpha\};\widetilde J)
  :=\widetilde{\mathcal M}^*_{0,T}
      (\widetilde M,\{\widetilde A_\alpha\};\widetilde J)/G_T.
\]
本节将在一点上偏离第~5 章的记号：我们总是假定\BookRecordIndex{T0340}稳定映射
$(\widetilde{\mathbf u},\mathbf z)\in
\widetilde{\mathcal M}^*_{0,T}
(\widetilde M,\{\widetilde A_\alpha\};\widetilde J)$ 的分支
$\widetilde u_0:S^2\to\widetilde M$ 满足
\[
  \pi\circ\widetilde u_0=\operatorname{id},
\]
因而是 $\widetilde M$ 的 $\widetilde J$-全纯截面。因此，\BookRecordIndex{T0413}重参数化群 $G_T$ 由元组
$(f,\{\phi_\alpha\}_{\alpha\in T})$ 组成，其中 $f:T\to T$ 是树自同构，并满足
$f(0)=0$、$\phi_0=\operatorname{id}$，而当 $\alpha\ne0$ 时
$\phi_\alpha\in G=\operatorname{PSL}(2,\mathbb C)$。它像定义~\ref{def:5-1-4} 中那样作用在
\BookRecordIndex{T0340}稳定映射\BookRecordIndex{T0246}模空间上。

现在考察这个\BookRecordIndex{T0246}模空间中某些\BookRecordIndex{T0066}标记点固定的子集。任取指标集
$I\subset\{1,\ldots,k\}$，并固定由 $I$ 标号的、$S^2$ 上两两不同的点所成的元组
\[
  \mathbf w=\{w_i\}_{i\in I}.
\]
这样的元组确定一个子集
\[
  \mathcal M^*_{0,T}
    (\widetilde M,\{\widetilde A_\alpha\};\mathbf w,\widetilde J)
  :=\widetilde{\mathcal M}^*_{0,T}
    (\widetilde M,\{\widetilde A_\alpha\};\mathbf w,\widetilde J)/G_T,
\]
它由满足下述条件的\BookRecordIndex{T0340}稳定映射等价类
$[\widetilde{\mathbf u},\mathbf z]\in
\mathcal M^*_{0,T}(\widetilde M,\{\widetilde A_\alpha\};\widetilde J)$ 组成：
\begin{equation}\label{eq:6-7-10}
  i\in I
  \quad\Longrightarrow\quad
  \pi\bigl(\widetilde u_{\alpha_i}(z_i)\bigr)=w_i.
\end{equation}
这意味着：当 $\alpha_i=0$，亦即 $z_i$ 位于特殊分支上时，$z_i=w_i$；否则，相应\BookRecordIndex{T0275}气泡
$\widetilde u_{\alpha_i}$ 的像包含在纤维 $\{w_i\}\times M$ 中。

可以按如下方式把这个\BookRecordIndex{T0246}模空间描述成\BookRecordIndex{T0283}求值映射的原像。如第~\ref{sec:6-2} 节，有一个单纯元组
$\widetilde{\mathbf u}=\{\widetilde u_\alpha\}_{\alpha\in T}$ 所成的\BookRecordIndex{T0246}模空间
$\mathcal M^*(\widetilde M,\{\widetilde A_\alpha\};\widetilde J)$，其中各
$\widetilde u_\alpha$ 都是表示类 $\widetilde A_\alpha$ 的 $\widetilde J$-全纯球面。只要
$\mathbf J\in\mathcal J_{\mathrm{reg}}(S^2;A_0)$，且对每个
$\alpha\in T\setminus\{0\}$ 都有
$\mathbf J\in\mathcal J_{\mathrm{reg}}^{\mathrm{Vert}}(S^2;A_\alpha)$，这个空间就是定向
有限维流形。仍如第~\ref{sec:6-2} 节，定义
\[
  Z(T)\subset(S^2)^E\times(S^2)^k
\]
为所有元组
$\mathbf z=(\{z_{\alpha\beta}\}_{\alpha E\beta},\{z_i\}_{1\le i\le k})$ 所成的集合；
对每个 $\alpha\in T$，其中满足 $\alpha E\beta$ 的点 $z_{\alpha\beta}\in S^2$，连同满足
$\alpha_i=\alpha$ 的点 $z_i\in S^2$，两两不同。于是有\BookRecordIndex{T0283}求值映射
\begin{equation}\label{eq:6-7-11}
  \widetilde{\operatorname{ev}}^E\times\pi^I:
  \mathcal M^*(\widetilde M,\{\widetilde A_\alpha\};\widetilde J)\times Z(T)
  \longrightarrow\widetilde M^E\times(S^2)^I,
\end{equation}
定义为
\[
  \widetilde{\operatorname{ev}}^E(\widetilde{\mathbf u},\mathbf z)
  :=\{\widetilde u_\alpha(z_{\alpha\beta})\}_{\alpha E\beta},
  \qquad
  \pi^I(\widetilde{\mathbf u},\mathbf z)
  :=\{\pi(\widetilde u_{\alpha_i}(z_i))\}_{i\in I}.
\]
满足式~\eqref{eq:6-7-10} 的单纯\BookRecordIndex{T0340}稳定映射\BookRecordIndex{T0246}模空间，是
$\widetilde{\operatorname{ev}}^E\times\pi^I$ 下子流形
$\widetilde\Delta^E\times\{\mathbf w\}$ 的原像。这引出下面的定义。

\begin{definition}\label{def:6-7-10}
若 $M$ 上被 $\omega$ \BookRecordIndex{T0372}驯服的\BookRecordIndex{T0088}殆复结构光滑族
$\mathbf J=\{J_z\}\in\mathcal J_\tau(S^2;M,\omega)$ 满足下列条件，就称它
\emph{关于 $\mathbf w$ 正则}。
\begin{description}[leftmargin=3.4em,style=nextline]
\item[\textup{(H)}] 对式~\eqref{eq:6-7-1} 的每个解 $u:S^2\to M$，$D_u$ 都满射。
\item[\textup{(V)}] 对每个 $z\in S^2$ 和每个单纯 $J_z$-全纯球面
$v:S^2\to M$，$D_{z,v}$ 都满射。
\sourcefootnote{原书的条件 (V) 遗漏了“单纯”。然而命题~\ref{prop:6-7-9} 只对
$v\in\mathcal M^*(A;J_z)$ 证明\BookRecordIndex{T0071}参数正则性，而定理~\ref{thm:6-7-11} 证明的步骤 4 又直接以
该命题\BookRecordIndex{T0326}推出 (V)；对任意多重覆盖球面要求同一结论一般过强，故依这两处上下文补正。}
\item[\textup{(E)}] 对每一棵带 $k$ 个标号、具有\BookRecordIndex{T0321}特殊顶点 $0$ 的树 $T$，以及每一族
$H_2(M;\mathbb Z)$ 中的同调类 $\{A_\alpha\}_{\alpha\in T}$，\BookRecordIndex{T0283}求值映射
~\eqref{eq:6-7-11} 都与 $\widetilde\Delta^E\times\{\mathbf w\}$ 横截。
\end{description}
关于 $\mathbf w$ 正则的\BookRecordIndex{T0088}殆复结构族之集记作
$\mathcal J_{\mathrm{reg}}(S^2;M,\omega;\mathbf w)$。
\end{definition}

准备至此，可以陈述如今已经标准的\BookRecordIndex{T0408}正则性定理。证明并不完全是例行公事，因为我们要求
\BookRecordIndex{T0088}殆复结构 $\widetilde J$ 保持 $\widetilde M$ 的纤维结构；见步骤 2。

\begin{theorem}\label{thm:6-7-11}
设 $\mathbf w=\{w_i\}_{i\in I}$ 是 $S^2$ 上两两不同的点所成的元组。
\begin{enumerate}[label=\textup{(\roman*)},leftmargin=3em]
\item 若 $\mathbf J\in\mathcal J_{\mathrm{reg}}(S^2;M,\omega;\mathbf w)$，则对每一棵具有
\BookRecordIndex{T0321}特殊顶点 $0$ 的树 $T$，以及每一族 $H_2(M;\mathbb Z)$ 中的球面同调类
$\{A_\alpha\}_{\alpha\in T}$，\BookRecordIndex{T0246}模空间
$\mathcal M^*_{0,T}
(\widetilde M,\{\widetilde A_\alpha\};\mathbf w,\widetilde J)$ 都是光滑定向流形，其维数为
\[
  \dim\mathcal M^*_{0,T}
  (\widetilde M,\{\widetilde A_\alpha\};\mathbf w,\widetilde J)
  =2n+2c_1(A)+2k-2|I|-2e(T),
\]
其中 $A:=\sum_{\alpha\in T}A_\alpha$。
\item $\mathcal J_{\mathrm{reg}}(S^2;M,\omega;\mathbf w)$ 是
$\mathcal J_\tau(S^2;M,\omega)$ 中的\BookRecordIndex{T0300}剩余集。
\end{enumerate}
\end{theorem}

\begin{proof}
由定义~\ref{def:6-7-10}，映射 $\widetilde{\operatorname{ev}}^E\times\pi^I$ 与
$\widetilde\Delta^E\times\{\mathbf w\}$ 横截，且
\[
  \widetilde{\mathcal M}^*_{0,T}
  (\widetilde M,\{\widetilde A_\alpha\};\mathbf w,\widetilde J)
  =(\widetilde{\operatorname{ev}}^E\times\pi^I)^{-1}
    (\widetilde\Delta^E\times\{\mathbf w\}).
\]
由命题~\ref{prop:6-7-7} 和命题~\ref{prop:6-7-9}，\BookRecordIndex{T0246}模空间
$\mathcal M^*(\widetilde M,\{\widetilde A_\alpha\};\widetilde J)$ 的维数为
\[
  \dim\mathcal M^*(\widetilde M,\{\widetilde A_\alpha\};\widetilde J)
  =2n+(2n+2)e(T)+2c_1(A).
\]
又因为
\[
  \dim Z(T)=4e(T)+2k,
  \qquad
  \operatorname{codim}(\widetilde\Delta^E\times\{\mathbf w\})
  =(2n+2)e(T)+2|I|,
\]
所以
\[
  \dim\widetilde{\mathcal M}^*_{0,T}
  (\widetilde M,\{\widetilde A_\alpha\};\mathbf w,\widetilde J)
  =2n+2c_1(A)+2k-2|I|+4e(T).
\]
群 $G_T$ 的维数为 $6e(T)$，且以保定向微分同胚作用在
$\widetilde{\mathcal M}^*_{0,T}
(\widetilde M,\{\widetilde A_\alpha\};\mathbf w,\widetilde J)$ 上，故 (i) 成立。

(ii) 的证明分四步进行；如第~\ref{sec:6-3} 节给出的命题~\ref{prop:6-2-8} 的证明那样，
对\BookRecordIndex{T0297}森林作归纳。步骤 2 最值得注意，因为我们在那里证明当前设置具有提供所需\BookRecordIndex{T0166}横截性的足够
灵活性。

\medskip
\noindent\textbf{步骤 1.}\quad
\BookRecordIndex{T0268}普适模空间
$\mathcal M^*(\widetilde M,\{\widetilde A_\alpha\};\mathcal J^\ell)$ 是\BookRecordIndex{T0053}巴拿赫流形，其中
$\mathcal J^\ell=\mathcal J_\tau^\ell(S^2;M,\omega)$ 表示 $M$ 上所有被 $\omega$ \BookRecordIndex{T0372}驯服的
$C^\ell$ \BookRecordIndex{T0088}殆复结构族所成的空间。

这是命题~\ref{prop:3-2-1} 与命题~\ref{prop:6-2-7} 证明中论证的改写。

\medskip
\noindent\textbf{步骤 2.}\quad
设 $T=(T,E,\Lambda)$ 是带 $k$ 个标号的\BookRecordIndex{T0297}森林，并设\BookRecordIndex{T0065}边求值映射
\[
  \widetilde{\operatorname{ev}}^{E,\mathcal J}:
  \mathcal M^*(\widetilde M,\{\widetilde A_\alpha\};\mathcal J^\ell)\times Z(T)
  \longrightarrow\widetilde M^E
\]
与 $\widetilde\Delta^E$ 横截。考察\BookRecordIndex{T0268}普适模空间
\[
  \widetilde{\mathcal M}^*_{0,T}
  (\widetilde M,\{\widetilde A_\alpha\};\mathcal J^\ell)
  :=(\widetilde{\operatorname{ev}}^{E,\mathcal J})^{-1}(\widetilde\Delta^E).
\]
则 $S^2$ 中两两不同的点所成的每个元组
$\mathbf w=\{w_i\}_{i\in I}$，都是普适投影
\[
  \pi^{I,\mathcal J}:
  \widetilde{\mathcal M}^*_{0,T}
  (\widetilde M,\{\widetilde A_\alpha\};\mathcal J^\ell)
  \longrightarrow(S^2)^I
\]
的\BookRecordIndex{T0409}正则值。

空间
$\widetilde{\mathcal M}^*_{0,T}
(\widetilde M,\{\widetilde A_\alpha\};\mathcal J^\ell)$ 由元组
\[
  (u_0,\{z_\alpha,u_\alpha\},\{z_{\alpha\beta}\},\{z_i\},\{J_z\})
\]
组成，其中 $z_\alpha$ 标记\BookRecordIndex{T0275}气泡 $u_\alpha$ 所在的纤维。这些元组满足
\begin{align*}
  \bar\partial_{\{J_z\}}(u_0)&=0,
  &\bar\partial_{J_{z_\alpha}}(u_\alpha)&=0
    &&(\alpha\in T\setminus\{0\}),\\
  0E\alpha&\Longrightarrow
  &u_0(z_{0\alpha})&=u_\alpha(z_{\alpha0}),
    &z_\alpha&=z_{0\alpha},\\
  \alpha E\beta,\quad\alpha\ne0,\quad\beta\ne0
  &\Longrightarrow
  &u_\alpha(z_{\alpha\beta})&=u_\beta(z_{\beta\alpha}),
    &z_\alpha&=z_\beta.
\end{align*}
为理解其切空间，注意：若 $S\subset T\setminus\{0\}$ 是\BookRecordIndex{T0417}子树，也就是说 $S$ 中任意两个
顶点都能由 $S$ 中的一条\BookRecordIndex{T0064}边链相连，则 $S$ 中所有\BookRecordIndex{T0275}气泡必定位于同一纤维内。因此，描述纤维
无穷小移动的向量 $\zeta_\alpha$ 对所有 $\alpha\in S$ 必须相同。此外，若某个顶点
$\beta\in S$ 由一条边与 $0$ 相连，则 $\zeta_\beta$ 必须等于描述根分支上附着点
$z_{0\beta}$ 无穷小移动的向量 $\zeta_{0\beta}$。所以，切空间由满足下列各式的元组
$(\xi_0,\{\zeta_\alpha,\xi_\alpha\},\{\zeta_{\alpha\beta}\},
\{\zeta_i\},\{Y_z\})$ 组成：
\begin{equation}\label{eq:6-7-12}
  D_{u_0}\xi_0+\frac12Y_z(u_0)du_0\circ j=0,
\end{equation}
\begin{equation}\label{eq:6-7-13}
  D_{u_\alpha}\xi_\alpha
  +\frac12Y_{z_\alpha}(u_\alpha)du_\alpha\circ j
  =\frac12\dot J_{\zeta_\alpha}(u_\alpha)du_\alpha\circ j
  \qquad(\alpha\ne0),
\end{equation}
\begin{equation}\label{eq:6-7-14}
  \xi_0(z_{0\alpha})+du_0(z_{0\alpha})\zeta_{0\alpha}
  =\xi_\alpha(z_{\alpha0})+du_\alpha(z_{\alpha0})\zeta_{\alpha0},
  \qquad \zeta_\alpha=\zeta_{0\alpha},
\end{equation}
\begin{equation}\label{eq:6-7-15}
  \xi_\alpha(z_{\alpha\beta})+du_\alpha(z_{\alpha\beta})\zeta_{\alpha\beta}
  =\xi_\beta(z_{\beta\alpha})+du_\beta(z_{\beta\alpha})\zeta_{\beta\alpha},
  \qquad \zeta_\alpha=\zeta_\beta.
\end{equation}
式~\eqref{eq:6-7-14} 对 $0E\alpha$ 成立，而式~\eqref{eq:6-7-15} 对满足
$\alpha E\beta$ 的 $\alpha,\beta\in T\setminus\{0\}$ 成立。式~\eqref{eq:6-7-13}
右端来自第 $\alpha$ 个\BookRecordIndex{T0275}气泡沿 $\zeta_\alpha$ 方向的水平移动。

$\pi^I$ 的微分为
\[
  (\xi_0,\{\zeta_\alpha,\xi_\alpha\},\{\zeta_{\alpha\beta}\},
    \{\zeta_i\},\{Y_z\})
  \longmapsto
  \left\{\begin{array}{ll}
    \zeta_i,&\alpha_i=0,\\
    \zeta_{\alpha_i},&\alpha_i\ne0
  \end{array}\right\}_{i\in I},
\]
其中照例由 $\alpha_i$ 标记包含\BookRecordIndex{T0066}标记点 $z_i$ 的分支。设相应基点
$(\widetilde{\mathbf u},\mathbf z)$ 满足
$\pi^I(\widetilde{\mathbf u},\mathbf z)=\mathbf w$，于是
\[
  w_i=\begin{cases}
    z_i,&\alpha_i=0,\\
    z_{\alpha_i},&\alpha_i\ne0,
  \end{cases}
  \qquad i\in I.
\]
\sourcefootnote{原书在这里误把微分方向的切向量元组代入非线性映射 $\pi^I$；但
$\pi^I$ 的定义域是\BookRecordIndex{T0246}模空间中的基点，紧接着的 $w_i=z_i$ 或 $z_{\alpha_i}$ 也只描述基点，
故改用相应的 $(\widetilde{\mathbf u},\mathbf z)$。}
任给切向量元组 $\widehat w_i\in T_{w_i}S^2$（$i\in I$），须证明存在满足式
~\eqref{eq:6-7-12}--\eqref{eq:6-7-15} 的元组
\[
  (\xi_0,\{\zeta_\alpha,\xi_\alpha\},\{\zeta_{\alpha\beta}\},
  \{\zeta_i\},\{Y_z\}),
\]
使得当 $i\in I$ 且 $\alpha_i=0$ 时
$\zeta_i=\widehat w_i$，而当 $i\in I$ 且 $\alpha_i\ne0$ 时
$\zeta_{\alpha_i}=\widehat w_i$。对 $\alpha_i=0$，这显然成立：只须取
$\zeta_i:=\widehat w_i$，并把元组的所有其他分量取为零。对 $\alpha_i\ne0$，论证如下。

取满足 $\alpha_i\ne0$ 的 $i\in I$，并考察相应\BookRecordIndex{T0417}子树
\[
  T_i:=\{\alpha\in T\setminus\{0\}\mid
    z_\alpha=w_i,\ [\alpha,\alpha_i]\ne\varnothing\}.
\]
连接顶点 $0$ 与 $T_i$ 的边至多有一条。（也可能一条都没有，因为这一步假定 $T$ 是\BookRecordIndex{T0297}森林，
而非树。）取
\[
  \zeta_{\alpha\beta}:=
  \begin{cases}
    \widehat w_i,&\alpha=0,\ \alpha E\beta,\ \beta\in T_i,\\
    0,&\text{其他情形},
  \end{cases}
  \qquad
  \zeta_\alpha:=
  \begin{cases}
    \widehat w_i,&\alpha\in T_i,\\
    0,&\alpha\in T\setminus(T_i\cup\{0\}),
  \end{cases}
\]
以及
\[
  \xi_\alpha:=0,
  \qquad \alpha\in T\setminus T_i.
\]
于是，式~\eqref{eq:6-7-15} 对
$\alpha,\beta\in T\setminus(T_i\cup\{0\})$ 成立，而式~\eqref{eq:6-7-14} 对满足
$0E\alpha$ 的 $\alpha\in T\setminus T_i$ 成立。因此，式~\eqref{eq:6-7-14} 和
~\eqref{eq:6-7-15} 等价于
\begin{equation}\label{eq:6-7-16}
  \xi_\beta(z_{\beta0})=du_0(w_i)\widehat w_i
  \quad(\beta\in T_i,\ 0E\beta),
  \qquad
  \xi_\alpha(z_{\alpha\beta})=\xi_\beta(z_{\beta\alpha})
  \quad(\alpha,\beta\in T_i,\ \alpha E\beta).
\end{equation}
\sourcefootnote{原书在式~\eqref{eq:6-7-16} 的第二个等式后只写
$\alpha,\beta\in T_i$，遗漏了 $\alpha E\beta$；\BookRecordIndex{T0191}节点 $z_{\alpha\beta}$ 仅沿树边定义，且该式
正是式~\eqref{eq:6-7-15} 沿边的化简，故补上边关系。}
所以，只须选取 $Y_z$ 与 $\{\xi_\alpha\}_{\alpha\in T_i}$，使式
~\eqref{eq:6-7-16} 成立，并满足
\begin{equation}\label{eq:6-7-17}
  Y_z(u_0(z))\equiv0,
  \qquad
  Y_{z_\alpha}(u_\alpha)\equiv0
  \quad\bigl(\alpha\in T\setminus(T_i\cup\{0\})\bigr),
\end{equation}
\begin{equation}\label{eq:6-7-18}
  D_{u_\alpha}\xi_\alpha
  +\frac12Y_{w_i}(u_\alpha)du_\alpha\circ j
  =\frac12\dot J_{\widehat w_i}(u_\alpha)du_\alpha\circ j
  \qquad(\alpha\in T_i).
\end{equation}
现在用引理~\ref{lem:3-4-3} 证明中的式~\eqref{eq:3-4-3} 来说明这些选择存在。回忆记号
$Z_\alpha:=\{z_{\alpha\beta}\mid\alpha E\beta\}$。

先假设没有边连接 $0$ 与 $T_i$。可以按如下方式进行。把 $Y_z$ 定义为
\begin{equation}\label{eq:6-7-19}
  Y_z:=\rho(z)\sum_{\alpha\in T_i}Y_\alpha,
\end{equation}
其中 $\rho:S^2\to[0,1]$ 是支集位于 $w_i$ 附近、且 $\rho(w_i)=1$ 的光滑\BookRecordIndex{T0188}截断函数；
$Y_\alpha$ 选取如下。若 $u_\alpha$ 是\BookRecordIndex{T0382}幽灵分支，即 $du_\alpha\equiv0$，就定义
$\xi_\alpha:=0$、$Y_\alpha:=0$。若 $du_\alpha\not\equiv0$，则用式
~\eqref{eq:3-4-3} 找到满足式~\eqref{eq:6-7-18} 的二元组
$(\xi_\alpha,Y_\alpha)$，使
$\xi_\alpha(Z_\alpha)=0$，对于满足 $z_\beta=w_i$ 的
$\beta\in T\setminus\{0,\alpha\}$ 有 $Y_\alpha(u_\beta)\equiv0$，并且
$Y_\alpha$ 在 $u_0(w_i)$ 附近消失。于是元组
$(\{Y_z\},\{\xi_\alpha\}_{\alpha\in T_i})$ 满足式~\eqref{eq:6-7-16}、
~\eqref{eq:6-7-17} 和~\eqref{eq:6-7-18}。

现在假设有一条边连接 $0$ 与 $T_i$，并以 $\beta_i\in T_i$ 表示满足 $0E\beta_i$ 的唯一
顶点。若 $\beta_i$ 是\BookRecordIndex{T0382}幽灵分支，以 $S_i\subset T_i$ 表示 $T_i$ 中所有能通过幽灵顶点链
与 $\beta_i$ 相连的幽灵顶点所成的\BookRecordIndex{T0417}子树。令 $V_i\subset T_i\setminus S_i$ 为所有不属于
$S_i$、但通过一条边与 $S_i$ 相连的顶点之集
\sourcefootnote{原书只说 $V_i\subset T_i$ 由“通过一条边与 $S_i$ 相连”的所有顶点组成，
未排除 $S_i$ 自身；按字面会纳入 $S_i$ 内部彼此相邻的幽灵顶点，并与紧接着的
“$u_\alpha$ 对每个 $\alpha\in V_i$ 都非常值”以及映射 $f$ 的唯一性冲突，故补明
$V_i\subset T_i\setminus S_i$。原书还把随后映射的靶写成 $S_i$，但在
$du_{\beta_i}\not\equiv0$ 的情形又令 $S_i=\varnothing$、$f(\beta_i)=0$；为同时涵盖这两种
情形，正文把靶补为 $S_i\cup\{0\}$。}，并以 $f:V_i\to S_i\cup\{0\}$ 表示满足
$\alpha E f(\alpha)$（$\alpha\in V_i$）的唯一映射。注意，对每个 $\alpha\in V_i$，
$u_\alpha$ 都非常值。若 $du_{\beta_i}\not\equiv0$，则定义
$S_i:=\varnothing$、$V_i:=\{\beta_i\}$，并记 $f(\beta_i):=0$。注意，对每个
$\alpha\in S_i$，都有 $u_\alpha\equiv u_0(w_i)$。定义
\[
  \xi_\alpha:=du_0(w_i)\widehat w_i,
  \qquad Y_\alpha:=0
  \qquad(\alpha\in S_i).
\]
于是式~\eqref{eq:6-7-18} 对每个 $\alpha\in S_i$ 成立。对 $\alpha\in V_i$，任选
$\xi'_\alpha\in W^{\ell,p}$，使
\[
  \xi'_\alpha(z_{\alpha f(\alpha)})=du_0(w_i)\widehat w_i,
  \qquad
  \xi'_\alpha\bigl(Z_\alpha\setminus\{z_{\alpha f(\alpha)}\}\bigr)=0.
\]
由式~\eqref{eq:3-4-3}，存在 $\xi''_\alpha\in W^{\ell,p}$ 与
$Y_\alpha\in C^\ell$，使
\[
  \xi''_\alpha(Z_\alpha)=0,
  \qquad
  D_{u_\alpha}(\xi'_\alpha+\xi''_\alpha)
  +\frac12Y_\alpha(u_\alpha)du_\alpha\circ j
  =\frac12\dot J_{\widehat w_i}(u_\alpha)du_\alpha\circ j;
\]
此外，$Y_\alpha$ 在 $u_0(w_i)$ 附近消失，并且对每个满足 $z_\beta=w_i$ 的
$\beta\in T\setminus\{0,\alpha\}$ 都有 $Y_\alpha(u_\beta)\equiv0$。定义
\[
  \xi_\alpha:=\xi'_\alpha+\xi''_\alpha.
\]

再考察其余顶点 $\alpha\in T_i\setminus(V_i\cup S_i)$。它们没有通过一条边与
$S_i\cup\{0\}$ 相连，因而可像第一种情形那样进行：若 $du_\alpha\equiv0$，就定义
$\xi_\alpha:=0$、$Y_\alpha:=0$；若 $du_\alpha\not\equiv0$，则用式
~\eqref{eq:3-4-3} 找到满足式~\eqref{eq:6-7-18} 的二元组
$(\xi_\alpha,Y_\alpha)$，使 $\xi_\alpha(Z_\alpha)=0$，对每个
$\beta\in T\setminus\{0,\alpha\}$ 且 $z_\beta=w_i$ 都有
$Y_\alpha(u_\beta)\equiv0$
\sourcefootnote{原书在这里仅要求 $\beta\in T_i\setminus\{\alpha\}$ 时消失；但\BookRecordIndex{T0297}森林中还可能
有不属于 $T_i$、却位于同一纤维 $w_i$ 的分支，而式~\eqref{eq:6-7-17} 要求这些分支不受
$Y_{w_i}$ 扰动。前面两次运用同一构造时也都对所有满足 $z_\beta=w_i$ 的其他分支施加了
消失条件，故据此补全。}，且 $Y_\alpha$ 在
$u_0(w_i)$ 附近消失。仍取支集位于 $w_i$ 附近的适当\BookRecordIndex{T0188}截断函数
$\rho:S^2\to[0,1]$，并用式~\eqref{eq:6-7-19} 定义 $Y_z$。于是
$(\{Y_z\},\{\xi_\alpha\}_{\alpha\in T_i})$ 如所需满足式~\eqref{eq:6-7-16}、
~\eqref{eq:6-7-17} 和~\eqref{eq:6-7-18}。

\medskip
\noindent\textbf{步骤 3.}\quad
对每个带 $k$ 个标号的\BookRecordIndex{T0297}森林 $T=(T,E,\Lambda)$，\BookRecordIndex{T0065}边求值映射
\[
  \widetilde{\operatorname{ev}}^{E,\mathcal J}:
  \mathcal M^*(\widetilde M,\{\widetilde A_\alpha\};\mathcal J^\ell)\times Z(T)
  \longrightarrow\widetilde M^E
\]
都与 $\widetilde\Delta^E$ 横截。

与命题~\ref{prop:6-2-8} 一样，可对\BookRecordIndex{T0297}森林归纳证明这一点：每一步把 $T'$ 中一对\BookRecordIndex{T0066}标记点
$(z_{\alpha\beta},z_{\beta\alpha})$ 变成新\BookRecordIndex{T0297}森林 $T$ 中一条边的两个端点。归纳步骤是证明
相应\BookRecordIndex{T0283}求值映射
\[
  \operatorname{ev}_{\alpha\beta}:
  \mathcal M^*_{0,T'}(\widetilde M,\{\widetilde A_\alpha\};\mathcal J^\ell)
  \longrightarrow\widetilde M\times\widetilde M
\]
与对角线横截。令
$\pi_{\alpha\beta}:=(\pi\times\pi)\circ\operatorname{ev}_{\alpha\beta}$。
由步骤 2，$\pi_{\alpha\beta}$ 横截于对角线
$\Delta_S\subset S^2\times S^2$
\sourcefootnote{原书在尚未定义 $\pi_{\alpha\beta}$ 时，称其与到 $S^2\times S^2$
的投影之复合映射横截；但复合的类型与紧接着的原像定义唯一确定
$\pi_{\alpha\beta}=(\pi\times\pi)\circ\operatorname{ev}_{\alpha\beta}$，原文亦把复合次序写反。}。
因此，由习题~\ref{ex:6-3-2}，只须再证明映射
\[
  \pi_M\circ\operatorname{ev}_{\alpha\beta}:
  \mathcal M^*_{0,T'}
    (\widetilde M,\{\widetilde A_\alpha\};\Delta_S,\mathcal J^\ell)
  :=\pi_{\alpha\beta}^{-1}(\Delta_S)
  \longrightarrow M\times M
\]
与对角线 $\Delta_M\subset M\times M$ 横截；这里
$\pi_M:\widetilde M\times\widetilde M\to M\times M$ 是自然投影。这一点来自定理
~\ref{thm:6-3-1} 证明中的论证。

\medskip
\noindent\textbf{步骤 4.}\quad
证明 (ii)。

由命题~\ref{prop:6-7-7} 和命题~\ref{prop:6-7-9}，满足定义~\ref{def:6-7-10} 中条件
(H) 与 (V) 的所有\BookRecordIndex{T0088}殆复结构
$\mathbf J\in\mathcal J_\tau(S^2;M,\omega)$，在
$\mathcal J_\tau(S^2;M,\omega)$ 中构成\BookRecordIndex{T0300}剩余集。现在令 $T$ 和
$\{A_\alpha\}_{\alpha\in T}$ 如定理所述。由步骤 2、3，还有一个\BookRecordIndex{T0268}普适模空间
\[
  \widetilde{\mathcal M}^*_{0,T}
    (\widetilde M,\{\widetilde A_\alpha\};\mathbf w,\mathcal J^\ell)
  :=(\pi^{I,\mathcal J})^{-1}(\mathbf w).
\]
考察投影
\[
  \pi_{T,\{A_\alpha\}}:
  \widetilde{\mathcal M}^*_{0,T}
    (\widetilde M,\{\widetilde A_\alpha\};\mathbf w,\mathcal J^\ell)
  \longrightarrow\mathcal J^\ell.
\]
这是弗雷德霍姆映射，并且 $\mathbf J\in\mathcal J^\ell$ 是
$\pi_{T,\{A_\alpha\}}$ 的\BookRecordIndex{T0409}正则值，当且仅当它对 $T$ 与 $\{A_\alpha\}$ 满足定义
~\ref{def:6-7-10} 中的条件 (E)。当 $\ell$ 充分大时，萨德--斯梅尔定理断言这些\BookRecordIndex{T0409}正则值在
$\mathcal J^\ell$ 中构成\BookRecordIndex{T0300}剩余集。从 $C^\infty$ 情形化归到 $C^\ell$ 情形，再次使用定理
~\ref{thm:3-1-6} 证明中的 Taubes 论证。因此，投影 $\pi_{T,\{A_\alpha\}}$ 的所有光滑
\BookRecordIndex{T0409}正则值在 $\mathcal J_\tau(S^2;M,\omega)$ 中构成\BookRecordIndex{T0300}剩余集。对所有 $T$ 与
$\{A_\alpha\}$ 取交，便知
$\mathcal J_{\mathrm{reg}}(S^2;M,\omega;\mathbf w)$ 是
$\mathcal J_\tau(S^2;M,\omega)$ 中的\BookRecordIndex{T0300}剩余集。所有这些细节都只是本书其他地方已有论证的
重复，留给读者在当前情形下验证。
\end{proof}

\begin{proof}[定理~\ref{thm:6-7-1} 的证明]
第一个断言立即来自定理~\ref{thm:6-7-11}\textup{(i)} 和命题~\ref{prop:6-1-2}。
只须再注意：\BookRecordIndex{T0246}模空间 $\mathcal M^*_{0,k}(A;\mathbf w,\mathbf J)$ 的紧化中新出现的每个边界
\BookRecordIndex{T0340}稳定映射，都以至少有一条边的树为模型
\sourcefootnote{原书说紧化中的“每个\BookRecordIndex{T0340}稳定映射”都有至少一条边；但紧化包含原主\BookRecordIndex{T0072}层，而主\BookRecordIndex{T0072}层
对应无边树。此处\BookRecordIndex{T0337}伪循环论证所考察的是逃离主\BookRecordIndex{T0246}模空间的序列所产生的新增边界极限，只有这些
\BookRecordIndex{T0340}稳定映射必有至少一条边，故据此限定。}，而且每个竖直 $J_z$-全纯球面都有非负 Chern 数。

为证明 (ii)，设对 $i=0,1$ 有
\[
  \mathbf J^i=\{J_z^i\}
  \in\mathcal J_{\mathrm{reg}}(S^2;M,\omega;\mathbf w^i)
    \cap\mathcal J_+(S^2;M,\omega;\kappa).
\]
若从 $(\mathbf w^0,\mathbf J^0)$ 到 $(\mathbf w^1,\mathbf J^1)$ 的同伦
$\{\mathbf w^\lambda,J_z^\lambda\}$ 满足以下条件，就称它正则。
\begin{description}[leftmargin=3.8em,style=nextline]
\item[\textup{(HV)}] 对每个 $A\in H_2(M;\mathbb Z)$，都有
\[
  \{J_z^\lambda\}
  \in\mathcal J_{\mathrm{reg}}(S^2;A;\mathbf J^0,\mathbf J^1)
  \cap\mathcal J_{\mathrm{reg}}^{\mathrm{Vert}}
    (S^2;A;\mathbf J^0,\mathbf J^1).
\]
\item[\textup{(E)}] 对每一棵带 $k$ 个标号、具有\BookRecordIndex{T0321}特殊顶点 $0$ 的树 $T$，以及每一族
$H_2(M;\mathbb Z)$ 中的同调类 $\{A_\alpha\}_{\alpha\in T}$，\BookRecordIndex{T0283}求值映射
\[
  \widetilde{\operatorname{ev}}^E\times\pi^I:
  \mathcal W^*(\widetilde M,\{\widetilde A_\alpha\};
    \{\widetilde J^\lambda\})\times Z(T)
  \longrightarrow\widetilde M^E\times[0,1]\times(S^2)^I
\]
与
$\widetilde\Delta^E\times
\{(\lambda,\mathbf w^\lambda)\}_\lambda$ 横截。这里，
$\mathcal W^*(\widetilde M,\{\widetilde A_\alpha\};
\{\widetilde J^\lambda\})$ 表示二元组 $(\lambda,\widetilde{\mathbf u})$ 所成的\BookRecordIndex{T0246}模空间，其中
$0\le\lambda\le1$，且
$\widetilde{\mathbf u}\in
\mathcal M^*(\widetilde M,\{\widetilde A_\alpha\};\widetilde J^\lambda)$。
\item[\textup{(P)}] 对每个 $\lambda\in[0,1]$，都有
\[
  \{J_z^\lambda\}_{z\in S^2}
  \in\mathcal J_+(S^2;M,\omega;\kappa).
\]
\end{description}
注意，条件 (E) 是有意义的，因为条件 (HV) 蕴含\BookRecordIndex{T0246}模空间
$\mathcal W^*(\widetilde M,\{\widetilde A_\alpha\};
\{\widetilde J^\lambda\})$ 是光滑有限维流形。从
$(\mathbf w^0,\mathbf J^0)$ 到 $(\mathbf w^1,\mathbf J^1)$ 存在\BookRecordIndex{T0407}正则同伦
$\{\mathbf w^\lambda,\mathbf J^\lambda\}$，这可由定理~\ref{thm:6-7-11} 证明中的类似论证
得到。这里还用到空间 $\mathcal J_+(S^2;M,\omega;\kappa)$ 是开且道路连通的。

与\BookRecordIndex{T0407}正则同伦 $\{\mathbf w^\lambda,J_z^\lambda\}$ 相伴的\BookRecordIndex{T0246}模空间
$\mathcal W^*_{0,k}(A;\{\mathbf w^\lambda,J_z^\lambda\})$，由元组
$(\lambda,u,\mathbf z)$ 组成，其中 $\lambda\in[0,1]$，
$u\in\mathcal M(A,S^2;\mathbf J^\lambda)$ 是式~\eqref{eq:6-7-1} 的解，并在该式中取
$J_z=J_z^\lambda$，且 $[u]=A$；而
$\mathbf z=(z_1,\ldots,z_k)$ 是 $S^2$ 中两两不同的点所成的元组，并在 $i\in I$ 时满足
$z_i=w_i^\lambda$。仅由条件 (HV) 和命题~\ref{prop:6-7-7} 已经可知，这个\BookRecordIndex{T0246}模空间是光滑
定向带边流形，维数为
\[
  \dim\mathcal W^*_{0,k}
  (A;\{\mathbf w^\lambda,J_z^\lambda\})
  =2n+2c_1(A)+2k-2|I|+1,
\]
且其边界为
\[
  \partial\mathcal W^*_{0,k}
  (A;\{\mathbf w^\lambda,J_z^\lambda\})
  =\mathcal M^*_{0,k}(A;\mathbf w^1,\{J_z^1\})
  \cup\bigl(-\mathcal M^*_{0,k}(A;\mathbf w^0,\{J_z^0\})\bigr).
\]
条件 (HV)、(E) 和 (P) 合起来说明，自然的\BookRecordIndex{T0283}求值映射
\begin{equation}\label{eq:6-7-20}
  \operatorname{ev}:\mathcal W^*_{0,k}(A;\{\mathbf w^\lambda,J_z^\lambda\})
  \longrightarrow M^k
\end{equation}
正是所需的\BookRecordIndex{T0337}伪循环配边。这就证明了 (ii)。

断言 (iii) 来自类似论证。唯一的区别是，此时 $\mathbf J^0$ 与 $z$ 无关，而
$\lambda=0$ 时的\BookRecordIndex{T0246}模空间只由单纯 $J^0$-全纯球面组成。以 $i\in I$ 为指标的三个\BookRecordIndex{T0066}标记点，
用来消去 $\lambda=0$ 时\BookRecordIndex{T0413}重参数化群 $G=\operatorname{PSL}(2,\mathbb C)$ 的作用。作此修改
以后，仍如上可知：对适当的同伦 $\{J_z^\lambda\}$，\BookRecordIndex{T0283}求值映射~\eqref{eq:6-7-20} 是\BookRecordIndex{T0337}伪循环
之间的配边。
\end{proof}

\begin{remark}\label{rem:6-7-12}
下面基于\BookRecordIndex{T0380}隐函数定理，概述定理~\ref{thm:6-7-1}\textup{(iii)} 的另一种证明。这个证明只能
说明两个\BookRecordIndex{T0337}伪循环（弱）表示同一个同调类。固定 $S^2$ 中两两不同的三个点所成的元组
\[
  \mathbf w=(w_1,w_2,w_3)
\]
以及\BookRecordIndex{T0088}殆复结构 $J\in\mathcal J_{\mathrm{reg}}(M,\omega)$。设
$X\subset M^k$ 是紧定向子流形，余维为
\[
  \operatorname{codim}X=2n+2c_1(A)+2k-6.
\]
可以选择 $X$，使得对每个 $T$ 和每个 $B$，\BookRecordIndex{T0283}求值映射
\[
  \operatorname{ev}_{J,T}:\mathcal M^*_{0,T}(B;J)\longrightarrow M^k
\]
都与 $X$ 横截。于是，\BookRecordIndex{T0283}求值映射
\[
  \operatorname{ev}_J:\mathcal M^*_{0,k}(A;\mathbf w,J)\longrightarrow M^k
\]
在\BookRecordIndex{T0337}伪循环的意义下与 $X$ \BookRecordIndex{T0281}强横截。现在取一列收敛到 $J$ 的
\[
  \{J_z^\nu\}\in
  \mathcal J_{\mathrm{reg}}(S^2;M,\omega;\mathbf w).
\]
由\BookRecordIndex{T0139}格罗莫夫紧性（定理~\ref{thm:5-5-5}），当 $\nu$ 充分大时，\BookRecordIndex{T0246}模空间
$\mathcal M^*_{0,k}(A;\mathbf w,\mathbf J^\nu)$ 的紧化中所有\BookRecordIndex{T0072}层上的\BookRecordIndex{T0283}求值映射都不交于
$X$。此外，由\BookRecordIndex{T0380}隐函数定理，$\mathcal M^*_{0,k}(A;\mathbf w,J)$ 的任意给定紧子集，在
$\nu$ 充分大时都嵌入 $\mathcal M^*_{0,k}(A;\mathbf w,\mathbf J^\nu)$。再用定理
~\ref{thm:5-3-1}，可以证明 $\operatorname{ev}_J:\mathcal M^*_{0,k}(A;\mathbf w,J)\to M^k$ 与 $X$
的交点，同 $\operatorname{ev}_{\mathbf J^\nu}:\mathcal M^*_{0,k}(A;\mathbf w,\mathbf J^\nu)\to M^k$
与 $X$ 的交点之间一一对应。因此，两个\BookRecordIndex{T0337}伪循环同每个子流形的交数都相同，故它们是同一个
同调类的弱表示。
\end{remark}

\paragraph{完整的\BookRecordIndex{T0145}格罗莫夫--威滕伪循环。}
最后，我们对更一般的\BookRecordIndex{T0143}格罗莫夫--威滕不变量作几点说明，并解释本书的结果如何纳入这一
背景。基本思路是考察\BookRecordIndex{T0283}求值映射与\BookRecordIndex{T0377}遗忘映射的乘积
\[
  \operatorname{ev}\times\pi:
  \overline{\mathcal M}_{0,k}(A;J)
  \longrightarrow M^k\times\overline{\mathcal M}_{0,k}.
\]
这个映射的靶空间是光滑紧流形（定理 D.4.2），而源空间是紧可度量空间
（定理~\ref{thm:5-6-6}）。暂且假装源空间也是流形，那么它带有基本类
$[\overline{\mathcal M}_{0,k}(A;J)]$；可以从
$M^k\times\overline{\mathcal M}_{0,k}$ 拉回上同调类，并在这个基本类上求值，以定义
\BookRecordIndex{T0143}格罗莫夫--威滕不变量。通过 Poincaré 对偶，可以如下解释这件事。取子流形
$W\subset\overline{\mathcal M}_{0,k}$，使它对偶于所要从
$\overline{\mathcal M}_{0,k}$ 拉回的类，并与 $\pi$ 横截。再考察其原像
\[
  \overline{\mathcal M}_{0,k}(A;W,J)
  :=\pi^{-1}(W)
  =\bigl\{(\mathbf u,\mathbf z)\in\overline{\mathcal M}_{0,k}(A;J)
      \mid\pi(\mathbf u,\mathbf z)\in W\bigr\},
\]
并在限制后的\BookRecordIndex{T0283}求值映射
\[
  \operatorname{ev}_W:\overline{\mathcal M}_{0,k}(A;W,J)\longrightarrow M^k
\]
下，对 $M^k$ 上上同调类的拉回作积分。习题~\ref{ex:6-6-5} 与定理~\ref{thm:6-7-1}
中的构造可以这样理解。这里必须取\BookRecordIndex{T0150}固定标记点
$\mathbf w=\{w_i\}_{i\in I}$ 所给出的子流形
$W\subset\overline{\mathcal M}_{0,k}$。用附录 D 的语言说，这相当于固定所有指标取自集合
$I$ 的\BookRecordIndex{T0185}交比。定理~\ref{thm:6-6-1} 的\BookRecordIndex{T0145}格罗莫夫--威滕伪循环对应于
$W=\overline{\mathcal M}_{0,k}$ 的情形，而\BookRecordIndex{T0150}固定标记点的\BookRecordIndex{T0283}求值映射则对应于
$W=\{\mathrm{pt}\}$ 的情形。

在能够证明\BookRecordIndex{T0246}模空间 $\overline{\mathcal M}_{0,k}(A;J)$ 确实带有自然基本循环
$[\overline{\mathcal M}_{0,k}(A;J)]^{\mathrm{virt}}$，且这个循环的同调类与 $J$ 无关之前，
以上仍只是启发式讨论。后一个条件的含义是：与
$\mathcal J_\tau(M,\omega)$ 中任意道路 $\{J_t\}_{0\le t\le1}$ 相伴的\BookRecordIndex{T0246}模空间，都带有一条
链，其边界是与 $J_0$ 和 $J_1$ 相伴的基本循环之差。若这些都成立，那么对
$\overline{\mathcal M}_{0,k}$ 的每个子流形 $W$，都可以定义同调类
\[
  \operatorname{ev}_*\bigl([\overline{\mathcal M}_{0,k}(A,[W];J)]^{\mathrm{virt}}\bigr)
  \in H_*(M^k),
\]
办法是对 $M^k$ 的每个子流形 $X$ 规定
\[
  \operatorname{ev}_*\bigl([\overline{\mathcal M}_{0,k}(A,[W];J)]^{\mathrm{virt}}\bigr)\cdot X
  :=(\operatorname{ev}\times\pi)_*
    [\overline{\mathcal M}_{0,k}(A;J)]^{\mathrm{virt}}\cdot[X\times W].
\]
这将在 $M^k$ 中给出只依赖于 $W$ 的同调类的良定义有理同调类。接下来的问题，是怎样对
$W$ 的特定选择，适当解释类
$[\overline{\mathcal M}_{0,k}(A,[W];J)]^{\mathrm{virt}}$。关于一般紧\BookRecordIndex{T0358}辛流形的
\BookRecordIndex{T0143}格罗莫夫--威滕不变量，近年来很多工作的目的正是如此；见
\srcite{127,239,249,343,376,189,281}。以与上述讨论相近的精神研究这一问题的工作，见
Cieliebak--Mohnke \,\srcite{68} 和 Ionel \,\srcite{197}。

如命题 7.4.3 所示，有一些特殊例子（带紧 Lie 群传递作用的 \BookRecordIndex{T0207}凯勒 流形），其中每个亏格
零\BookRecordIndex{T0340}稳定映射都正则。利用复结构的可积性，可以证明在这种情形下，\BookRecordIndex{T0246}模空间
$\overline{\mathcal M}_{0,k}(A;J)$ 自然具有\BookRecordIndex{T0135}复轨形结构。（见 Fulton--Pandharipande
\srcite{133}、Behrend--Manin \,\srcite{34}、Robbin--Ruan--Salamon \,\srcite{333}。）
所以，在这种情形下，$\overline{\mathcal M}_{0,k}(A;J)$ 确实带有自然有理基本循环。
不过，在这个基本循环上的积分是否真是辛结构的形变不变量，远非显然。

本书并不试图解释紧化 $\overline{\mathcal M}_{0,k}(A;J)$ 的同调。我们所证明的是：在
\BookRecordIndex{T0058}半正情形，对通有\BookRecordIndex{T0088}殆复结构，\BookRecordIndex{T0283}求值映射在最高\BookRecordIndex{T0072}层上的限制由定理~\ref{thm:6-6-1} 是\BookRecordIndex{T0337}伪循环。
同样，由习题~\ref{ex:6-6-3}，它在每一\BookRecordIndex{T0072}层
$\mathcal M^*_{0,T}(\{A_\alpha\};J)$ 上的限制都是\BookRecordIndex{T0337}伪循环。再把\BookRecordIndex{T0088}殆复结构的类别扩展到依赖于
$z\in S^2$ 的族 $\{J_z\}$，定理~\ref{thm:6-7-1} 说明：\BookRecordIndex{T0283}求值映射在最高\BookRecordIndex{T0072}层的子集
$\mathcal M^*_{0,k}(A;\mathbf w,\{J_z\})$ 上的限制也是\BookRecordIndex{T0337}伪循环；在该子集上，点
$w_i$（$i\in I$）固定。

来看看这在 $k=4$ 的简单情形下意味着什么。$\overline{\mathcal M}_{0,4}$ 的最高\BookRecordIndex{T0072}层可与
$S^2\setminus\{0,1,\infty\}$ 等同，而三个非平凡\BookRecordIndex{T0072}层各由一点组成。因此，取
$[W]=[\mathrm{pt}]$，就有两种不同方法定义应与
$[\overline{\mathcal M}_{0,4}(A,[\mathrm{pt}];J)]^{\mathrm{virt}}$ 对应的循环。第一种是
\[
  \operatorname{ev}_T:\mathcal M^*_{0,T}(\{A_\alpha\};J)\longrightarrow M^4,
\]
其中 $T$ 是一个非平凡\BookRecordIndex{T0072}层；第二种是
\[
  \operatorname{ev}:\mathcal M^*_{0,4}(A;\mathbf w,\{J_z\})\longrightarrow M^4,
\]
其中 $I=\{1,2,3,4\}$，而 $\mathbf w$ 是
$\overline{\mathcal M}_{0,4}$ 最高\BookRecordIndex{T0072}层中的一个固定点。从这些定义完全看不出两个循环表示
同一个同调类：前者由两条相交曲线表示，每条曲线各带两个\BookRecordIndex{T0066}标记点；后者则由带四个\BookRecordIndex{T0066}标记点、
\BookRecordIndex{T0185}交比固定的单条曲线表示。这两个循环同调，是一个深刻结果，也是\BookRecordIndex{T0256}粘合定理的一个特例。我们
将在定理 7.5.10 中以更一般的形式陈述它，并在第~10 章证明。第~7.2 节算出了一个简单
例子；见注 7.2.10。

自然会问，乘积
\[
  \operatorname{ev}\times\pi:\mathcal M^*_{0,k}(A;J)
  \longrightarrow M^k\times\overline{\mathcal M}_{0,k}
\]
是否也是\BookRecordIndex{T0337}伪循环。由习题~\ref{ex:6-5-13} 和引理~\ref{lem:6-5-14}，这意味着对通有元素
$\mathbf w\in\overline{\mathcal M}_{0,k}$，
\[
  \operatorname{ev}_{\mathbf w}:\mathcal M^*(A;J)\longrightarrow M^k
\]
是\BookRecordIndex{T0337}伪循环。例~\ref{ex:6-6-4} 已经讨论了怎样尝试证明这一点。那里把 $J$ 限制为固定且不
依赖于 $z$ 的\BookRecordIndex{T0088}殆复结构，却遇到如下困难：命题~\ref{prop:6-1-2} 中到某个底\BookRecordIndex{T0072}层单纯\BookRecordIndex{T0340}稳定映射
的约化，一般无法选成保持\BookRecordIndex{T0066}标记点的\BookRecordIndex{T0185}交比。定理~\ref{thm:6-7-1} 已经表明，只要允许 $J$
依赖于 $z\in S^2$，就能绕过这一问题。这足以说明每个 $\operatorname{ev}_{\mathbf w}$ 都是\BookRecordIndex{T0337}伪循环。
不过，$\operatorname{ev}\times\pi$ 是\BookRecordIndex{T0337}伪循环这一陈述还要更强，它还要求控制
$\operatorname{im}(\operatorname{ev}\times\pi)$ 中那些位于 $\overline{\mathcal M}_{0,k}$ 所有\BookRecordIndex{T0072}层上方的
极限点。对一般\BookRecordIndex{T0058}半正 $(M,\omega)$ 做到这一点，需要考察方程更为一般的扰动。

我们在这里不继续讨论，因为这个一般构造不会产生新的不变量：\BookRecordIndex{T0339}稳定曲线\BookRecordIndex{T0246}模空间
$\overline{\mathcal M}_{0,k}$ 的同调由与\BookRecordIndex{T0084}带标号树相伴的子流形生成（定理 D.6.4），而由
习题~\ref{ex:6-6-3}，我们已经知道它们给出\BookRecordIndex{T0337}伪循环。不过，值得指出的是，若对
$(M,\omega)$ 作相当强的限制性假设，则与习题~\ref{ex:6-6-5} 类似的论证确实可以说明，
相应的映射到 $M^k\times\overline{\mathcal M}_{0,k}$ 是\BookRecordIndex{T0337}伪循环。此外，在这种情形，借助
命题~\ref{prop:6-5-17} 中的构造，可以自然解释相应的\BookRecordIndex{T0337}伪循环
$\operatorname{ev}_{\mathbf w}$。下面的习题请读者验证细节。

\begin{exercise}\label{ex:6-7-13}
设 $(M,\omega)$ 是最小 Chern 数为 $N$ 的单调\BookRecordIndex{T0358}辛流形，并设
$J\in\mathcal J_{\mathrm{reg}}(M,\omega)$。固定同调类 $A\in H_2(M;\mathbb Z)$、整数
$k\ge3$ 和指标集 $I\subset\{1,\ldots,k\}$，并满足
\[
  3\le\#I\le N+2.
\]
以 $\overline{\mathcal M}_{0,I}$ 表示\BookRecordIndex{T0066}标记点指标集为 $I$ 的（光滑）亏格零\BookRecordIndex{T0339}稳定曲线\BookRecordIndex{T0246}模空间
（见附录 D）。
\begin{enumerate}[label=\textup{(\roman*)},leftmargin=3em]
\item 证明乘积映射
\[
  \operatorname{ev}\times\pi_I:\mathcal M^*_{0,k}(A;J)
  \longrightarrow M^k\times\overline{\mathcal M}_{0,I}
\]
是\BookRecordIndex{T0337}伪循环，其中 $\pi_I$ 是删去映射 $\mathbf u$ 以及\BookRecordIndex{T0066}标记点
$z_i$（$i\notin I$）的\BookRecordIndex{T0377}遗忘映射。

\emph{提示：}先考察 $I=\{1,\ldots,k\}$ 的情形。$\operatorname{ev}\times\pi$ 的极限集由满足
\[
  \mathbf w=\lim_{\nu\to\infty}\pi([\mathbf z^\nu]),
  \qquad
  x_i=\lim_{\nu\to\infty}u^\nu(z_i^\nu),
  \quad 1\le i\le k,
\]
的点
$(\mathbf x,\mathbf w)\in M^k\times\overline{\mathcal M}_{0,k}$ 组成；这里序列
$(u^\nu,\mathbf z^\nu)\in\mathcal M^*_{0,k}(A;J)$ 没有极限位于
$\mathcal M^*_{0,k}(A;J)$ 中的收敛子列。现在分析\BookRecordIndex{T0066}标记点在极限中的行为。这个论证与习题
~\ref{ex:6-6-5} 中的十分相似，只是极限\BookRecordIndex{T0339}稳定曲线 $\mathbf w$ 的结构比先前更复杂。

\item 设
$\mathbf w\in\mathcal M_{0,I}\subset\overline{\mathcal M}_{0,I}$ 是开\BookRecordIndex{T0072}层中的一点。证明
$\operatorname{ev}\times\pi_I$ 与 $M^k\times\{\mathbf w\}$ \BookRecordIndex{T0294}弱横截。把习题~\ref{ex:6-6-5} 的\BookRecordIndex{T0337}伪循环
$\operatorname{ev}_{\mathbf w}$ 解释为\BookRecordIndex{T0283}求值映射在
$\mathcal M^*_{0,k}(A;\mathbf w,J)=\pi_I^{-1}(\mathbf w)$ 上的限制；这是
$\mathcal M^*_{0,k}(A;J)$ 的一个子流形（见命题~\ref{prop:6-5-17}）。
\end{enumerate}
\end{exercise}
